\section{The component estimate on the selected levels}
\label{sec:flag-component-transfer}

The degree calculation gives a number for each component of a deep cover of the
original core. The division estimate, however, uses a lattice at a selected level. To
pass between them we compare component measures and follow the lattice through the
level maps. The two signed
refinements retain their own levels. This transfer provides the
component estimate used in the common comparison of
Section~\ref{subsec:joint-family-estimate}. The component index
\(Z\) always refers to the deep cover in
Lemma~\ref{hyp:covolume-deep-component}. For each CM label \(\tau\)
the determinant line used in the transfer is the ratio of the two
Lie-eigenspace determinants. Definition~\ref{hyp:flat-determinant-frames}
specifies its frame metrics, and Proposition~\ref{prop:deep-flat-line-degree}
computes the normalized degree.

\begin{established}[Nonterminal component and level]\label{hyp:covolume-deep-component}
At the nonterminal lift \(j\) of the original core, put
\((p,q)=(p_j,q_j)\) and write \(N_{\rm lev}\) for the deep level.
The signature, Hermitian rank and level are
\[
 (p_j,q_j)=(k+j-1,m-j),\qquad r=p_j+q_j=k+m-1,\qquad
 N_{\rm lev}=g_0g_1\,15^e .
\]
Here \(0<q<p\), \(f(r-1)>2\), and both marked type primes divide
\(N_{\rm lev}\).
\end{established}
\begin{proof}
The ordered row in Section~\ref{subsec:paired-row-indices} moves one
negative coordinate to the positive block at each lift. Its initial
sizes are \(k-1,m\), and \(j<m\) gives \(q>0\). The selected
ranges \(k\ge m^4\), \(m\ge 2\) and \(f\ge8\) give the strict
inequalities. Section~\ref{subsec:deep-cover-on-core} takes principal
level \(g_0g_1\,15^e\) and auxiliary level one, so
\[
 \{\ell:\ell\mid N_{\rm lev}\}\subseteq\{3,5,g_0,g_1\}.
\]
\end{proof}

\begin{lemma}[Equality of component measures]\label{lem:deep-component-measures}
For the nonterminal lift and level of
Lemma~\ref{hyp:covolume-deep-component}, write \(w_c\) for the component weights in the normalized flag degree and \(\mathcal C\) for the component set. The weights are equal, so
\[
 w_c=\frac1{|\mathcal C|}\sum_{a\in\mathcal C}w_a\qquad(c\in\mathcal C).
\]

\end{lemma}

\begin{proof}
\proofstep{Congruence groups in a common Hermitian space.}
Complex uniformization at a fixed embedding identifies each
weight \(w_c\), up to a common factor, with the Riemannian
covolume of an arithmetic subgroup on the product of special
unitary domains. Transport the full-level polarized lattices
of two components to the same Hermitian space with its CM type.
The arithmetic groups are the inverse images of congruence
stabilizers on the finite adeles. An integral unitary
determinant has all its conjugates of modulus one and hence
is torsion.

Let \(n\ge3\) be the principal level and \(u\) the determinant.
If \(u\ne1\), its congruence and the bound at every complex embedding
would give
\[
 n^{[E:\Q]}\le |N_{E/\Q}(u-1)|
       =\prod_{\sigma:E\hookrightarrow\C}|\sigma(u)-1|
       \le 2^{[E:\Q]} .
\]
Thus the arithmetic group at this level has determinant one.
Write \(F=E^+\), \(H=\operatorname{SU}(V)\) for the common
Hermitian space, and let \(K_i\subset H(\mathbb A_{F,f})\) be
the two finite level stabilizers. In the moduli uniformization
the components lie in the same finite adelic orbit of marked
lattices. A finite adelic unitary isometry \(g\) therefore gives
\[
 K_2=gK_1g^{-1},\qquad K_{12}=K_1\cap K_2,\qquad
 \Gamma_i=H(F)\cap K_i\quad(i=1,2,12).
\]
\label{cor:covolume-component-transfer-1}

\proofstep{The index ratio and strong approximation.}
The group \(H\) is simply connected and absolutely almost
simple. Since \(p,q>0\) at a real place, the archimedean product
has a noncompact factor. Strong approximation thus gives
\[
         \overline{H(F)}=H(\mathbb A_{F,f})
\]
in the finite adelic topology (Platonov--Rapinchuk,
\emph{Algebraic Groups and Number Theory}, Theorem~7.12).
\cite[the strong approximation theorem]{milne-pel-period-domain}

Each coset of \(K_{12}\) in \(K_i\) is open
and meets \(H(F)\), hence
\[
 [\Gamma_i:\Gamma_{12}]=[K_i:K_{12}]
    =\frac{\operatorname{vol}_f(K_i)}{\operatorname{vol}_f(K_{12})}
      \qquad(i=1,2).
\]
Conjugation by the finite adelic unitary element \(g\) preserves
Haar measure on \(H\). Indeed, over a splitting field its adjoint
action on \(\mathfrak{sl}_r\) has determinant one. Conjugation on
all matrices has determinant \((\det g)^r(\det g^{-1})^r=1\)
and fixes the complementary scalar line. The same determinant
identity holds at every finite place. The two indices are therefore
equal. Both real quotients are covered by the quotient for
\(\Gamma_{12}\), of degrees \([\Gamma_i:\Gamma_{12}]\).
Dividing its volume by these equal degrees proves \(w_c=w_{c'}\).
\label{cor:covolume-component-transfer-2}

\end{proof}

\paragraph{Passage from intersections to counts.}
To turn the component degree into a count, start with the integral strictly decreasing
weight \(u\) with
gap parameter \(g\) and \(B=B_0(u,g)\) of the preceding
section. For a moduli component \(Z\) let \(c\) be its flag
component under the component bijection. Write
\(W=w_cG(u)>0\), \(a(s)\) for the arithmetic
self-intersection on \(c\), and \(d(s)\) for its geometric
degree. We have \(d(s)\to W\) and \(a(s)/(nd(s))\to s_c\)
along positive rationals \(s\downarrow0\). Let \(\delta\)
denote the normalized top degree from cusp-form ranks on \(Z\).
Write \(V_Z(t)\) for its logarithmic count volume at exponential
radius \(e^{kt}\), on multiples \(k\) of eight, normalized by the
field degree and arithmetic Hilbert growth factor.

\begin{lemma}[Volume slope on a component]
Put \(c_0=2h_{\rm HT}(B)\|u\|_1\) and \(s'_Z=s_c-c_0\). Then
\[
 V_Z(t)\ge W(s'_Z+t)_+\ge\delta(s'_Z+t)_+\qquad(t\in\mathbb R).
\]
For \(E(u,g)=\ell_B\|u\|_1+q\sum_\tau U_\tau\) the arithmetic
mean satisfies \(\operatorname{mean}_Z s'_Z\ge-(C_*+88)E(u,g)\).
\end{lemma}
\begin{proof}
At the staircase weight \(u^0\) of successive gaps one the nef
limiting sections dominate cusp growth, with degree \(w_cG(u^0)\).
The normalized cusp degree at \(u\) multiplies the reference growth
by \(G(u)/G(u^0)\), hence \(\delta\le w_cG(u)=W\). \label{cor:covolume-component-transfer-3}
The arithmetic section-volume inequality for the perturbed
semipositive, fibrewise nef line, whose sections are cuspidal, gives
for some finite \(A\) depending on the model and weight, and every
sufficiently small positive rational \(s\),
\[
 V_Z(t)\ \ge\
 \max\{0,\ a(s)/n+d(s)(t-c_0-As)\},\qquad
 c_0=2h_{\rm HT}(B)\|u\|_1 .
\]
Here \(h_{\rm HT}(B)\) is the normalized finite Hodge--Tate cost.
The archimedean comparison shrinks the logarithmic radius by \(O(s)\).

At a finite place it costs \(s a_v+2 h_v(B)\|u\|_1\) with
finite support, where \(\sum_v h_v(B)/[F:\Q]=h_{\rm HT}(B)\).
After division by \([F:\Q]\), their total is
\[
 s\,\frac{\sum_v a_v}{[F:\Q]}
       +2\|u\|_1\frac{\sum_v h_v(B)}{[F:\Q]}
   = O_{Z,u}(s)+c_0 .
\] These comparisons use the
actual integral mixed metrics of \(Z\).
\label{cor:covolume-component-transfer-4}

Since \(d(s)\to W\), \(a(s)/(nd(s))\to s_c\), and \(A\) is fixed
on the model, the expression inside the positive part tends to
\(W(s_c+t-c_0)\). This proves the pointwise estimate. By
\eqref{eq:flag-slope-lower} the mean of \(s_c\) is at least
\(-C_*E(u,g)\); the finite-place sum gives
\(c_0\le88\ell_B\|u\|_1\le88E(u,g)\). Subtraction gives the mean
bound. \label{cor:covolume-component-transfer-5}
\end{proof}

\paragraph{Field of definition.}
The component bounds will be combined over a coefficient field. To keep track of this
field, observe that the toroidal system in the slope computation can be chosen on
the moduli space of the given deep cover. At each residue
characteristic a factor of level at least three survives
removal of that characteristic's power, since one of
\(g_0,g_1\) survives. Enlarge the coefficient field \(F\)
so as to contain both the required roots of unity and a field
of definition of the coefficient curve at the auxiliary level.
The CM field, ranks and signatures are preserved under this
change of the datum. To descend the inequality to \(F\),
use the bijection of marked components preserving the
idempotents in their evaluation rules. Cusp growth has the
same top degree after extension, by expansion in a field basis.
\label{cor:covolume-component-transfer-6}

\begin{lemma}[Descent of the count inequality]
Let \(F'/F\) be a finite extension over which the coefficient curve
and the required level roots of unity are defined, with the marked
components identified under scalar extension. Write \(\mu_{F'},\mu_F\) for the corresponding normalized component slopes. Every lower bound descends,
\[
 \mu_{F'}\ge B\quad\Longrightarrow\quad\mu_F\ge B.
\]

\end{lemma}
\begin{proof}
Write \(d=[F':F]\), and take an \(F\)-basis \(b_i\) of \(F'\)
with trace dual \(b_i^\vee\). Choose \(N_0>0\) making all
\(N_0b_i^\vee\) integral. On \(f=\sum_i b_i f_i\) the maps
\(f\mapsto\operatorname{Tr}_{F'/F}(N_0b_i^\vee f)=N_0f_i\)
are jointly injective. Trace preserves the component and cusp
conditions. On each fixed mixed chart evaluation commutes with trace;
integrality of the scaled dual elements thus gives the required
finite integrality tests for every image.

At each
embedding, summing the absolute values of the conjugates of
\(N_0b_i^\vee\) bounds the norm of its image by \(A'\|f\|_\infty\),
with \(A'\) independent of \(k\). For the integral norm balls
in the two fields this gives
\[
 \#B_{F'}(e^{kt})\le\#B_F(A'e^{kt})^d .
\]
For any \(\eta>0\), use this inequality at \(t-\eta\) and take
\(k\) large enough that \(A'e^{-k\eta}\le1\). The factor \(d\)
in the logarithm cancels against the field degree. Thus the lower
bound at \(t-\eta\) over \(F'\) is also a lower bound at \(t\)
over \(F\); let \(\eta\downarrow0\) in its continuous positive part.
\label{cor:covolume-component-transfer-7}
\end{proof}

\begin{auxlemma}[A thick-frame lower bound for the evaluation norm]
\label{lem:deep-thick-evaluation}
Fix a nonterminal quotient weight \(u_1\) on a deep component \(Z\).
Write \(\mathcal B_{u_1,Z}(y)\) for the squared operator norm of
evaluation at the unit flag frame \(y\), on the cusp space with
the Petersson measure normalized to mass one on \(Z\). The same lower
bound holds for the restriction to \(Z\) of the probability measure
on the full marked object, since its component mass is at most one. Let \(\lambda\) denote the first
minimum of the polarized trace lattice in the positive majorant at
that frame, and let \(M_Z\) be its principal congruence level. There is
an absolute \(C\) such that, on the selected family,
\[
 M_Z\lambda^2\ge b
  \quad\Longrightarrow\quad
 \mathcal B_{u_1,Z}(y)\ge e^{-B_1},\qquad
 B_1=C\{\|u_1\|_1+(fr)^8\}.
\]
The constant is independent of the component, lift and coefficient embedding.
\end{auxlemma}
\begin{proof}
Pass to an orthonormal flag frame by an upper triangular change. The normalized point norm, the positive majorant and the Petersson integral are invariant under this change. In the lower triangular coordinates \(I+z\) used in Lemma~\ref{lem:short-weight-curvature}, set
\[
       d=f\binom r2,\qquad R=(100r^2)^{-1}.
\]
If a group element \(\gamma\) identifies two points of \(\|z\|\le R\), it increases majorant lengths by at most a factor two. Indeed the operator norm of \(z\) at each infinite place is at most \(R\); on each initial coordinate subspace the Gram matrix therefore differs from the diagonal sign matrix by at most \(2R+R^2\). Successive orthogonal projections in triangular elimination change the form by a factor at most \((1-4rR)^{-r}<\sqrt2\). Comparing each of the two flag frames with the center gives the claimed factor.

In particular the coordinate change satisfies
\[
           \|\gamma v\|_y\le 2\|v\|_y\qquad(v\in\Lambda\otimes\mathbb R).
\]

The polarized integral alternating pairing identifies the trace lattice with its Euclidean dual via the orthogonal complex structure of the majorant. In real dimension \(2fr\) transference consequently supplies a rational basis of lattice vectors of length at most \((2fr)^2/\lambda\). If \(\gamma\ne1\), it moves some member \(v\) of this basis. Principal congruence and the length comparison give
\[
      M_Z\lambda\le\|(\gamma-1)v\|
          \le3\|v\|\le3(2fr)^2/\lambda .
\]
This contradicts \(M_Z\lambda^2\ge b> 12(fr)^2\) for large scale. Thus this chart injects on the component. \cite{banaszczyk}

We construct the evaluation section on the complex flag quotient. Use the invariant Kähler metric given by a strictly positive flag weight, with gaps one inside the two blocks and a cross-block gap of order \(r\). It is complete: a homogeneous Riemannian metric on the flag domain is complete and its quotient by a properly discontinuous group is complete. In a unitary root coordinate \(z_{ij}\) the Levi form of a weight \(v\), using the sign convention of Lemma~\ref{lem:short-weight-curvature}, has coefficient
\[
           (v_i-v_j)J_j/J_i\qquad(i<j).
\]

The normalized line of the section is tensored with the inverse canonical bundle when applying the top-degree \(L^2\) estimate. This subtracts the volume weight, whose successive differences have modulus two. The divided row has consecutive gaps at least \(T\) before passing to the flag order. In that order the within-block gaps and the opposite signed gap across the two blocks have the positive signs in the displayed formula. Unitary invariance gives this curvature comparison on the entire quotient.

Put \(\rho=R/2\), choose a fixed smooth \(0\le\chi\le1\) equal to one for \(t\le1/4\) and zero for \(t\ge1\), and use the local holomorphic triangular frame of norm one at its center to extend \(F_0=\chi(\|z\|^2/\rho^2)\) by zero. At its center introduce the nonpositive weight
\[
 \psi_\delta(z)=d\chi(\|z\|^2/\rho^2)
          \log\big((\|z\|^2/\rho^2+\delta)/2\big),\quad 0<\delta\le1 .
\]

The triangular frame is holomorphic throughout the ball, so
\[
   \operatorname{supp}\bar\partial F_0
       \subset\{\rho/2\le\|z\|\le\rho\}.
\]

On the inner ball the Levi form of \(\psi_\delta\) is nonnegative. On the transition annulus the modulus of the logarithm is bounded by three and its first two radial derivatives are bounded by absolute constants in the variable \(\|z\|^2/\rho^2\). Thus throughout the chart
\[
  i\partial\bar\partial\psi_\delta
       \ge-Cd\rho^{-2}i\partial\bar\partial\|z\|^2
       \qquad(0<\delta\le1).
\]
 The triangular change to an orthonormal chart and its inverse have bounded derivative on this ball: the inverse Gram matrices on initial subspaces have norm at most \((1-2R-R^2)^{-1}\) and triangular elimination gives the bound \((1-4rR)^{-2r}\). The positive gaps absorb this loss since
\[
          d\rho^{-2}\ll fr^6,\qquad T\ge (frL)^{24}.
\]

Solve \(\bar\partial u_\delta=\bar\partial F_0\) with this weight by the complete Kähler \(L^2\) estimate. In the unit frame the invariant metric of the weight line is a product of powers of successive Gram determinants. On the ball every successive pivot has absolute value between \(1/2\) and \(2\), so both this metric and its inverse are at most \(\exp(C\|u_1\|_1)\). 
The volume density has the analogous bound \(\exp(Cdr)\). On the support of the derivative, \(\psi_\delta\ge-3d\) and \(\|\bar\partial F_0\|\ll\rho^{-1}\) in the chart norm. Integration over the chart and replacement of its volume by the larger component volume give, in probability measure,
\[
       \int_Z\|u_\delta\|^2e^{-\psi_\delta}
       \le \exp\{C(\|u_1\|_1+(fr)^8)\}.
\]

Weak compactness on an exhaustion as \(\delta\downarrow0\) gives a solution \(u\) with the limiting bound. It is holomorphic near the center. If \(u(0)\ne0\), the limiting norm is infinite since its radial integral contains
\[
           \int_0^{\rho/2} t^{-2d} t^{2d-1}\,dt=\infty .
\]
 Hence \(s=F_0-u\) is holomorphic, \(s(y)=1\), and its squared Petersson norm satisfies the same exponential bound after enlarging \(C\). \cite[Theorem~5.1]{demailly-complete-kahler} \cite{hormander-slice}

For these regular weights the holomorphic square-integrable section is cuspidal. At a rational boundary use the Fourier expansion in the unipotent radical. Integration over its period parallelepiped bounds the norm below by the contribution of any one Fourier coefficient. If that coefficient is nonzero, its continuous point norm is bounded below on a compact transverse coordinate ball of positive measure in the boundary chart. For an index \(\nu\) vanishing on a rational boundary ray \(Y\), the factor along that ray satisfies
\[
    \left|e^{2\pi i\langle\nu,x+itY\rangle}\right|
                         = e^{-2\pi t\langle\nu,Y\rangle}=1.
\]
Thus this Fourier mode cannot gain radial decay. The normalized squared frame factor grows with exponent at least the positive cross-block gap, whereas the Haar density loses at most \(2fr^2\) in that radial coordinate (the sum of positive root multiplicities with their radial weights). Writing \(t\ge0\) for logarithmic radial length, such a coefficient would contribute at least a positive multiple of
\[
             \int_0^\infty\exp[(T-2fr^2)t]\,dt=\infty .
\]

Negative frequencies likewise violate integrability. This is the zero-cutoff case of the cusp-flow calculation in the proof of Lemma~\ref{lem:short-weight-curvature} and the ensuing Fourier estimate. The analytic identification of the specified marked holomorphic row in Lemma~\ref{lem:marked-full-holomorphic-row}, applied with this deep congruence marking, realizes the section in the cusp space at level \(M_Z\). The quotient \(|s(y)|^2/\|s\|_{\rm P}^2\) gives the assertion.
\end{proof}

\begin{auxlemma}[Evaluation at a frame with a short period]
\label{lem:deep-thin-evaluation}
Let \(Z\) be a component at a nonterminal lift, \(u_1\) one of its
selected quotient weights and \(y\) a frame on \(Z\) at a good
complex embedding. Write \(N_*=g_0g_1 15^e\) for the deep level,
\(\lambda=\lambda_1(y)\) for the first minimum in the polarization
metric and \(\mathcal B_{u_1,Z}(y)\) for the supremum of
\(\|g(y)\|^2/\|g\|_{2,Z}^2\) over nonzero cusp sections of weight
\(u_1\) on \(Z\), with component measure of mass at most one. Uniformly
over these components, weights and frames, for large family scale,
\[
 N_*\lambda^2<b\quad\Longrightarrow\quad
 \mathcal B_{u_1,Z}(y)\ge e^{-\phi(y)/625},\qquad
 \phi(y)=\frac{2f}{\operatorname{rd}(E)^2\lambda^2}.
\]
\end{auxlemma}
\begin{proof}
\proofstep{Choose a differential operator at the quotient weight.}
Put \(M_0=g_0g_1\), \(d_*=2fm\), \(R=\operatorname{rd}(E)\).
Use the specific column count \eqref{eq:deep-quotient-column-count} and exponents \eqref{eq:deep-quotient-exponents} for \(u_1\). The
corresponding theta weight \(G\) is obtained by the parallel shift \(w_\tau=n/2+1-u_{\tau,p}\) and central exponent \(n/2\). If an embedding
has \(q\) rows with exponents \(a_1\ge\cdots\ge a_q\), use one
\(j\)-row minor \(P_{\tau j}\) at each height. By the definition of
\(n\), all their column sets may be chosen disjoint. Put
\[
 P=\prod_{\tau}\prod_{j=1}^{q_\tau}
          P_{\tau j}^{\,a_{\tau j}-a_{\tau,j+1}},
 \qquad a_{\tau,q_\tau+1}=0 .
\]

This polynomial is harmonic since no two contracted opposite labels occur in
the same column. Put \(D=\deg P\). The positive last exponent gives
cuspidality: at a rational boundary the projected frame loses rank
in the corresponding full minor, so that its constant term vanishes,
as in Lemma~\ref{lem:marked-class-section-cuspidal}. A determinant
of size \(j\) contributes at most \(j!\le r^j\) to the sum of
absolute coefficients. Hence
\[
 1\le D\le fr(4b+1)\le b^2,\qquad \|P\|_1\le r^D .
\]

The host, level and column estimates in
\ref{cor:deep-boundary-orders-5} and
Lemma~\ref{lem:joint-theta-parameters} give, uniformly,
\[
 frn\le b,\quad R^7\le b,\quad d_*\le2b,
 \qquad b^{16}\le M_0\le b^{70}.
\]

Indeed \(f,r,n\) grow polynomially in \(x\), whereas
\(b=\exp(x^\epsilon)\); at each type prime
\((kb)^8\le g_i\), and the upper type bounds give
\(\log g_i\le35\log b\). In the joint level estimate use this
value of \(D\) and the smaller signature bound \(q_{\rm small}\le m\);
the numerical condition is
\[
 R_0^2R^{28}Dr=90(fr+D)R^{28}Dr=O(b^9)\le M_0 .
\]
The inequalities \(m^4\le f\) and \(R\ge\max(16,f^2)\) are independent of the weight.

\proofstep{Compare one actual translate with all frames on the original level.}
For the frame over \(y\) after forgetting the deep marking, use
the same lattice and form. Its first minimum is still \(\lambda\).
Choose the marked half \(\mathfrak m=(\bar\alpha)\), \(\alpha\bar\alpha=M_0\).
The isometry of Lemma~\ref{lem:scaled-light-metric} sends \(v\mapsto\bar\alpha v\) from the original lattice in the metric \(H_A\) to the half-level lattice in \(H'=H_A/M_0\). It preserves the first minimum, since
\[
 H'(\bar\alpha v,\bar\alpha v)=H_A(v,v).
\]

The isogeny pushes Haar probability forward to Haar probability, and the marked letter has squared length \(2\xi_R/M_0\) in \(H'\). Thus all \(M_0^{frn}\) translations of the
principal theta tensor are readings of cusp rules at weight \(G\).
Write \(b_\nu(z)\) for their invariant point norms multiplied by
the common Lie normalization \((2\pi\xi_R)^{-D/2}\).
Apply Lemmas~\ref{lem:light-thick-lift}, \ref{lem:light-frequency-column}
and~\ref{lem:light-block-jet-product} to this particular \(P\).
These lemmas treat arbitrary such ordered disjoint blocks. Changing to the
marked derivative costs \((2\xi_R/M_0)^{D/2}\). Parseval over the
half-level translates, followed by the maximum bound for a finite sum,
gives
\[
 \begin{aligned}
 M_0^D\sum_\nu b_\nu(y)^2
    &\ge \pi^{-D}[24(2d_*)^{d_*/2}]^{-n}
            2^{-(r+2)D}e^{-2a\phi(y)},\\
 a&=8100\pi n m^4d_*^4R^7/M_0,\\
 -\log\max_\nu b_\nu(y)&\le
    3b^3D+a\phi(y)+(D+frn)\log M_0 .
 \end{aligned}
\]

For the last inequality, \(D\ge1\), \(d_*\le2b\) and \(r,n\le b\)
bound half the logarithm of all powers of \(2,\pi,24,2d_*\) in
the first line by \(3b^3D\). Also
\[
 a\le 8100\pi\,16 b^{10}/b^{16}\le 1/2500
 \qquad(b\ge1000).
\]
This gives a nonzero rule at \(y\); fix a maximizing translate \(\nu\).

For its upper bound on the original level, let \(z_0\) be an orthonormal frame with value at least half its supremum. The rank-one
period estimate in \ref{cor:slice-short-heights-3}, applied with
spread at most \(3b^2\), gives first minimum at least
\(e^{-40L}\). Thus, with \(L=3\log b\),
\[
 2b_\nu(z_0)\ge\sup_z b_\nu(z),\qquad \lambda_1(z_0)\ge e^{-40L}.
\]

The four bounds needed here are
\[
 \log M_0\le30L,\qquad
 \log(fr)\le L,\qquad
 \log R\le3L,\qquad
 D+frn\le2b^2.
\]

Set \(h=frn\) and \(A=2f^{3/2}n^2R\).
For the standard tensor lattice a rational basis in its second
factor consists of integral vectors of lengths at most \(A\). Contraction with this basis
bounds the tensor minimum below by \(e^{-40L}/A\).
Use the submean and Cauchy inequalities in
\ref{cor:slice-short-heights-4} for the translated theta function
in dimension \(h\). A polydisc of coordinate radius
\[
 b_0=\min\{1/\sqrt{\pi h},e^{-40L}/(4A\sqrt h)\}
\]
embeds in the torus.

Its submean estimate costs
\(e^{1/2}(\sqrt\pi b_0)^{-h}\). A multi-index of degree \(D\)
costs \((\pi eD)^{D/2}\) times the coefficient sum by Cauchy's
inequality on the Gaussian function. The Lie factor cancels the
change from unit to marked coordinates. Including the factor two
for the half-supremum frame gives
\[
 \log\sup_z b_\nu(z)\le \log 2+\tfrac12
      +h\log(1/(\sqrt\pi b_0))+D\log r+\tfrac D2\log(eD)
       =O(b^2\log b).
\]
Here \(h\le b\), \(D\le b^2\), and \(\log A=O(\log b)\).
One may also retain the full finite marking count
\(M_0^{frn}\) in the supremum estimate since its logarithm has
the same bound.

Since \(e=\lfloor3\log b/\log3\rfloor\), eventually \(15^e\ge b\).
The thin-frame condition gives \(M_0\lambda^2\le1\) and hence
\[
 \phi(y)\ge b^{14},\qquad
 3b^3D+(D+frn)\log M_0+O(b^2\log b)\le \phi(y)/2500 .
\]
Together with the bound on \(a\), this proves
\(e^{-\phi(y)/1250}\sup_z b_\nu(z)\le b_\nu(y)\).

\proofstep{Return to the selected weight on the given component.}
Forget the refined marking and pull the chosen cusp rule back to
the deep level. In the adelic description the map enlarges the
level stabilizer on the right; the Hermitian lattice and the
orthonormal frame are carried by the same rational isometry.
The first minimum and every invariant point norm are preserved.
Write \(\pi\) for the forgetting map and \(e_Z\) for the idempotent of \(Z\). For the selected translate \(f_\nu\),
\[
 \|e_Z\pi^* f_\nu(z)\|=\|f_\nu(\pi z)\|\quad(z\in Z),
 \qquad (e_Z\pi^* f_\nu)|_{Z^c}=0 .
\]

The parallel determinant generator of Lemma~\ref{lem:flat-generic-generator},
raised at each embedding to power \(-w_\tau\), and the inverse central unit raised to power \(n/2\) in the weight return of
\eqref{eq:parallel-return-norm} and \ref{cor:deep-cover-on-core-10} convert the weight \(G\) to
\(u_1\). Both factors are invertible on the component and have
constant positive invariant norm there. Their product multiplies
both sides of the ratio estimate by one and the same positive
constant. The resulting section \(g\) is cuspidal and
\[
 \|g(y)\|\ge e^{-\phi(y)/1250}\|g\|_{\sup,Z}>0,\qquad
 \|g\|_{2,Z}\le\|g\|_{\sup,Z}.
\]
Squaring proves the assertion for \(\mathcal B_{u_1,Z}(y)\).
\end{proof}

\paragraph{Use of the shifted weight.}
The intersection estimate applies at a shifted weight. To recover the height needed on the quotient of weight \(u_1\), take its prescribed
shift \(u'=u'(u_1)\), the gap \(g\) in its budget and the
bounds \(s'_Z\) just proved for \(u'\). Write \(D_Z>0\)
for the image degree at \(u_1\) and \(H_Z\) for its normalized Chow
height in Proposition~\ref{prop:component-arithmetic-degree}. A hyperplane
has the Petersson norm at infinity and the component section lattice
at finite places. Write \(c_s\) for a log bound for its short integral
generators, and \(c_m\ge0\) for the norm cost of the carrying
section from \(u'\). Take \(h=h_{\rm HT}(B)\|u'\|_1\).
\begin{established}[Rank shares and length costs on the selected quotient]\label{hyp:covolume-shift-costs}
Put \(J_Z=\binom{T-1}{d}D_Z\), \(z=7(d^2/T+d\rho)\) and
\(U=fmbL\log\log x\). Uniformly for the nonterminal quotient weights and the deep components, beyond a common scale,
\[
  (1-z)r_Z\le J_Z\le r_Z,\qquad
  T\max(E(u',g),h)\le C_b U,
\]
\[
    [6d+10+T(z+2\rho)]c_s\le C_sU,\qquad Tc_m\le C_mU.
\]
Here \(u'=u_1-3r\omega_{\rm H}\); the carrying cost is measured per unit weight and includes the finite places.
\end{established}
\begin{proof}
The moving rank and the upper bound at the same quotient weight are proved in Propositions~\ref{prop:deep-moving-rank-lower} and~\ref{prop:deep-shifted-rank-upper}. They give \(J_Z\le r_Z\le(1+z)J_Z\); use \((1-z)(1+z)\le1\). The logarithm of the generator norm satisfies \(c_s\le P(\kappa/T+1)\) by the component generator bound in the Hilbert--Samuel calculation. Since \(d\ge2\), \(z\ge7d^2/T\) and \(z\ge7d\rho\), we have
\[
 6d+10+T(z+2\rho)\le4Tz,
 \qquad 4Tz c_s\le4zP\kappa+4zTP\le8b .
\]
The last step uses \(zP\kappa\le b\), \(z\le1/2\) and \(TP\le\sqrt b\), the inequalities used in \eqref{eq:deep-hs-budget}.

For the arithmetic cost let \(S\) bound the head width before division and let \(c_0\) be the determinant shift there. The \(k-1\) head entries are centered by their common reference entry; only the \(m=r-k+1\) tail entries retain a displacement of order \(b\). Euclidean division at length \(T\), including the staircase and the Hodge correction, therefore gives
\[
 T\|u'\|_1\le f\{rS+8mb+r(c_0+4r)\}\le834fmb .
\]

Indeed the selected reserve gives \(P S\le100w''\le800b\), \(r\le P\) and \(c_0\le22 f^2r^4\), while \(f^2r^5\le P\le b\). The same bound applies before the Hodge correction. Write \(h_0\) for the undivided split-gap parameter. The quotient satisfies
\[
   \|u'\|_\infty\le 8b/T,\qquad g=h_0/(4T),\qquad
   fr^2\le P\le h_0,\quad
   1\le B=4fr^2(1+\|u'\|_\infty/g)\le132b.
\]

Here the positive quotients at the smaller block have distance at least \(g\) from all other entries at their label. For completeness write \(w_\tau\) for the undivided tail scale and \(I_\tau\) for the \(m=r-k+1\) outer indices. The divided-array estimates read
\[
 U_\tau\le 3w_\tau/(2T),\qquad
       |u'_{\tau i}|\ge w_\tau/(4T)\quad(i\in I_\tau).
\]

Indeed dividing the extreme entries costs at most a unit per coordinate. The staircase and Hodge adjustment together cost at most \(c_0+4r\) before division; the separation reserve in the selected quotient absorbs this and the head width into \(w_\tau/4\). Since \(q\le m\),
\[
        q U_\tau\le 6\sum_{i\in I_\tau}|u'_{\tau i}|.
\]
 From \(\log B\le L+\log132\), the period identity \(|\log\Omega|\le3L\), the discriminant bound and the prime allowance \(\{3,5,g_0,g_1\}\) one obtains
\[
 \ell_B+6\le100L,\qquad h_{\rm HT}(B)\le100L.
\]
For the finite term sum \(\log p/(p-1)\) over the rational primes up to \(B\) by the factorial bound \(\sum_{p\le B}\log p/(p-1)\ll1+\log B\); at the primes of the level there are at most four additional summands. These estimates prove the two length bounds for \(E\) and \(h\), after one enlargement of \(C_b\).

It remains to pay for the section taking weight \(u'\) to \(u_1\). Use the same product of \(r\) theta null values at labels generating the marked quotient modulo \(5^4\) as in Lemma~\ref{lem:deep-boundary-uniform-order}, choosing one product \(F\) nonzero on \(Z\). Put \(a(y)=\|F(y)^2\|\) in the normalized metric and \(\phi=2f/(\operatorname{rd}(E)^2\lambda^2)\).
The single-tuple estimate \eqref{eq:eta-period-decay} yields a \(P_0\ll(frL)^2\) with
\[
 \log a(y)\le P_0/625\quad\text{everywhere},\qquad
 \log a(y)\le (P_0-\phi)/625\quad(M_Z\lambda^2<b).
\]

For the second comparison we need the active part of the torus bound, write \(M_0=g_0g_1\). The deep level satisfies \(M_Z=M_0 15^e\ge bM_0\); hence \(M_Z\lambda^2<b\) implies \(\lambda<M_0^{-1/2}\le h\) for the torus cutoff. In \eqref{eq:eta-period-decay} put \(R=5\), square the product and absorb its logarithmic power in the difference between \(\pi\phi/625\) and \(\phi/625\). The resulting constant is bounded by the single-tuple calculation in Proposition~\ref{prop:auxiliary-scalar-section}.

Divide the evaluation norm estimate according to the same two ranges. Lemmas~\ref{lem:deep-thick-evaluation} and~\ref{lem:deep-thin-evaluation} give, with \(\beta=P_0+B_1\),
\[
        a(y)\le e^\beta,\qquad
        a(y)^3\le e^\beta\mathcal B_{u_1,Z}(y),
        \qquad T\beta\ll U .
\]
For the length bound, \(TP\le\sqrt b\), \(P\ge(frL)^{12}\) and \(T\|u_1\|_1\le834fmb\) absorb both polynomial terms in \(P_0+B_1\). The exponent \(-3\phi/625\) in the thin range is sufficient against \(-\phi/625\) in the lower bound for the evaluation norm.

The carrying section is \(F^{24}\) times the inverse central root to power \(3r\) and the flat line of parallel weight \(-12r\). The central degrees cancel, and the resulting weight is \(24r\omega_{\rm H}=8(u_1-u')\). The finite orders of the two constant factors enter the degree of that flat section; the remaining theta factor has the boundary order of Lemma~\ref{lem:deep-boundary-uniform-order}. In the normalized metric the identity for the eighth multiple reads
\[
      \operatorname{wt}(\theta)=8(u_1-u'),\qquad \|F^{24}\|=a^{12}.
\]
 To read the contribution per unit of the shift, raise the Bergman estimate to the eighth power and take the sixteenth root:
\[
      a^{24}\le e^{8\beta}\mathcal B^8,\qquad
      \|F^{24}\|^{1/8}=a^{3/2}\le e^{\beta/2}\mathcal B^{1/2}.
\]
The logarithmic norm cost is therefore at most \(\beta/2\). Summing the orders in the flat line and using the determinant degree computation of Proposition~\ref{prop:deep-flat-line-degree} gives the remaining increment
\[
 \frac{3r\cdot2fr\log(2\pi)}{8}
       -\frac{12r\cdot f(p-q)\log\Omega}{16}
       \ \le\ 10P .
\]

Indeed \(|p-q|\le r\), \(|\log\Omega|\le3L\), and \(P\ge(frL)^{12}\). The scalar used to represent the flat line cancels by the product formula. Thus \(c_m\le \max(0,\beta/2+10P)\) and \(Tc_m\ll U\). This completes the four uniform estimates.
\end{proof}
\begin{lemma}[The shifted slope at the selected length]
Suppose the shifted degree satisfies
\(\delta'_Z\ge(1-\rho)D_Z\), \(0\le\rho<1\), on the deep good model.
Use the component-degree bound \eqref{eq:deep-component-degree}, the
Hilbert--Samuel error budget \eqref{eq:deep-hs-budget}, and the
rank shares and length costs in Lemma~\ref{hyp:covolume-shift-costs}.
Write \(\mu_{\rm deep}(u_M)\) for the squared-metric slope of the deep cusp lattice at \(u_M=u_W-c_0\omega_{\rm H}\). Then
\[
 \mu_{\rm deep}(u_M)\ge-C_U U,\qquad C_U=2(3C_1+C_2+4).
\]

\end{lemma}
\begin{proof}
Short generation gives \(H_Z+c_s\ge0\). Multiplication by the
carrying section and the birational upper count for the image give
for every \(\eta>0\)
\[
 \delta'_Z(s'_Z+c_s)\le V_Z(c_s)
 \le D_Z(H_Z+c_s+c_m+h+\eta).
\]
For \(s'_Z+c_s>0\) divide by \(D_Z\) and use the degree
comparison. For \(s'_Z+c_s\le0\) the right side is nonnegative.
Letting \(\eta\) decrease to zero gives the component inequality
\[
 (1-\rho)\max(0,s'_Z+c_s)\le H_Z+c_s+c_m+h .
\]

In the length budget for each quotient the length times
\(E(u',g)\) and the length times \(h\) are bounded by
\(C_b U\), where \(U\) is the error unit on this run.
The short-generator and carrying costs have the corresponding
bounds with constants \(C_s,C_m\). In particular the two length costs in the shifted height are
\[
 T(c_m+h)\le(C_m+C_b)U,\qquad
 T(C_*+88)E(u',g)\le(C_*+88)C_b U.
\]
Together with the short generators these give
\(C_1=C_s+C_m+C_b+(C_*+88)C_b\).
\label{cor:covolume-component-transfer-8}

Here the short-generator cost has its own multiplier
\(6d+10+T(z+2\rho)\), where \(d\) is the flag dimension,
\(T\ge3d+5\) is the chosen length and \(z<1\) the share
loss. We give the step that turns the component height into
a bound for the full lattice. Put
\(J_Z=\binom{T-1}{d} D_Z\) and let \(r_Z\) be the section
rank in the unshifted weight.

For \(0\le z<1\), Lemma~\ref{hyp:covolume-shift-costs} gives
\[
                    (1-z)r_Z\le J_Z\le r_Z .
\]
Equation~\eqref{eq:deep-component-degree} puts all \(D_Z\) in
\([D_*,D_* /(1-\rho)]\) for some \(D_*>0\).
Set \(X=(C_*+88)E(u',g)\), \(c=c_m+h\) and
\(H_Z^+=H_Z+c_s\). The preceding component inequalities imply
\[
 \frac{\sum_Z J_Z H_Z^+}{\sum_Z r_Z}
 \ge (1-z)(1-\rho)^2(c_s-X)-c .
\]
Indeed \(\sum r_Z\le \#\{Z\}\min J_Z/((1-z)(1-\rho))\) and
\(\operatorname{mean}_Z(s'_Z+c_s)_+\ge c_s-X\).

In the numerator
bound the positive terms using \(\min J_Z\); for the term
\(-c\sum J_Z\) use \(\sum J_Z\le\sum r_Z\).
Equation~\eqref{eq:deep-hs-budget} bounds the ordinary
Petersson degree of the deep geometric cusp lattice of
weight \(u_M=u_W-c_0\omega_{\rm H}\) below by
\[
 J_Z\{T H_Z-3(d+1)|H_Z|-A\log(2+|H_Z|)-R\},
 \qquad A\log\max(A,1)+R\le C_2 U ,
\]
with \(A>0\) and \(R\ge0\).
Here \(C_2=10240\) suffices since \(b\le U\).
The logarithmic term can be absorbed into the linear error:
\[
 A\log(2+|H_Z|)\le A\log\max(A,1)+2+|H_Z| .
\]
Indeed, apply \(\log v\le v-1\) to \(v=(2+|H_Z|)/\max(A,1)\).
Use \(|H_Z|\le H_Z^++c_s\) in the bracket. Its coefficient
of \(H_Z^+\) is at least \(T-3d-4>0\).

Put \(b_d=3d+4\) and \(\alpha=(1-z)(1-\rho)^2\).
After summation and division by \(\sum_Z r_Z\), the coefficient
lost on \(c_s\) is bounded by
\[
 (T+b_d)-(T-b_d)\alpha
      = T(1-\alpha)+b_d(1+\alpha)
      \le T(z+2\rho)+6d+8 .
\]
Here \(0\le\alpha\le1\), \(1-\alpha\le z+2\rho\) and
\(\sum J_Z\le\sum r_Z\) bound the negative contributions.
The other losses are at most \(TX+Tc+C_2U+2\).
Each of the short, \(TX\) and \(Tc\) terms is at most
\(C_1U\). The squared-norm slope is at least twice the sum
of component degrees divided by the sum of their ranks.
Thus for \(U\ge1\) the squared-metric lower bound on the
deep quotient lattice at \(u_M\) holds
with \(C_U=2(3C_1+C_2+4)\). \label{cor:covolume-component-transfer-9}

\end{proof}

We have a slope bound at the shifted weight. Before it can be used on the row we must
restore its determinant weight and account for the change in metric. Let
\(\Lambda_M\) be the deep cusp lattice at \(u_M\) with
its squared Petersson metric and let \(\Lambda_-\) be the
deep row lattice at the negative shift. The flat line
from \(\Lambda_M\) to \(\Lambda_-\) restores the determinant weight of the row
and the prescribed theta root. Let
\(w_\tau\) be the integral parallel shifts at the embedding labels, \(\omega\) the
logarithmic CM period in the specified data, and \(c_{1,\tau}\) the
flat-line coefficients in the slope comparison.

\begin{paramdefinition}[Component, determinant line and frame metrics]\label{hyp:flat-determinant-frames}
Fix a component \(Z\) of the deep cover over its coefficient field \(F\),
and a label \(\tau\) of the CM field \(E\). Write \(V_\tau,V_{\bar\tau}\)
for the two Lie eigenspaces, of dimensions \(p,q\). For frames \(v,w\) put
\[
 \mathcal L_\tau=\det V_\tau\otimes(\det V_{\bar\tau})^{-1},\qquad
 e(v,w)=(v_1\wedge\cdots\wedge v_p)\otimes
                      (w_1\wedge\cdots\wedge w_q)^{-1}.
\]
The flat weight is \((-1,+1)\) on the two determinants, with zero
theta exponent. Indeed, if \(\phi=f(v,w)e(v,w)\), then
\[
 f(vA,wB)=\frac{\det B}{\det A}f(v,w).
\]
Use the determinant metric induced by the polarization. At a finite
place take the integral eigenframes of the given prime-to-residue
level model. Write \(\mathfrak d_E\) for the different,
\(R_E=\log\operatorname{rd}(E)\), and \(D_0'=\log d_0\), where \(d_0\)
is the positive discriminant parameter of the marked quadratic field.
For the odd quadratic character \(\chi\) of \(\mathbb Q(\sqrt{-d_0})\),
use the CM normalization
\[
              \omega=\frac{L'(0,\chi)}{L(0,\chi)}+D_0'-R_E .
\]
\end{paramdefinition}

\begin{construction}
\begin{tikzpicture}[>=Stealth,every node/.style={align=center,inner sep=2pt}]
\node (line) at (0,0) {\(\mathcal L_\tau\) with generator \(\phi\)};
\node (arch) at (-2.75,-1.12) {Archimedean frames\\\(\|\phi\|_\sigma=c_\sigma\)};
\node (fin) at (2.75,-1.12) {Integral frames over \(w\)\\constant normalized order};
\node (cm) at (0,-2.35) {Compare on a CM frame};
\node (deg) at (0,-3.48) {\(\widehat{\deg}_{\rm flat}(\mathcal L_\tau)
        =\tfrac{p-q}{2}\omega+q(D'_0-R_E)\)};
\draw[->] (line) -- (arch);
\draw[->] (line) -- (fin);
\draw[->] (arch) -- (cm);
\draw[->] (fin) -- (cm);
\draw[->] (cm) -- (deg) node[midway,right=1mm] {\(\scriptstyle-\sum_w\delta_{w,\tau}\log{\rm N}w/[F:\mathbb Q]\)};
\end{tikzpicture}
\caption{Archimedean norms and finite orders of the flat generator in Definition~\ref{hyp:flat-determinant-frames}. The finite correction converts the CM frame degree to the normalized degree of Proposition~\ref{prop:deep-flat-line-degree}.}\label{fig:flat-cm-frames}
\end{construction}

\begin{proposition}[The normalized degree of the flat line]
\label{prop:deep-flat-line-degree}
For each component \(Z\) and label \(\tau\) of
Definition~\ref{hyp:flat-determinant-frames}, the determinant rule gives a
generator \(\phi\) on the proper normal generic fibre, up to scalar.
Its norm is constant on the archimedean frame tests at each embedding. At a finite place its normalized order is constant on the integral
frame tests. With the finite-order correction
\[
 \delta_{w,\tau}=q\{\operatorname{ord}_w(\tau\mathfrak d_E)
                              -\operatorname{ord}_w(d_0)\},
\]
the normalized degree is
\[
 \widehat{\deg}_{\rm flat}(\mathcal L_\tau)
       =\frac{p-q}{2}\,\omega+q(D_0'-R_E).
\]
Thus the contribution to the slope coefficient, with the squared-norm convention of this comparison, is
\[
 -\widehat{\deg}_{\rm flat}(\mathcal L_\tau)=-\tfrac{p-q}{2}\omega+q(R_E-D'_0).
\]

\end{proposition}

\begin{proof}
Choose the generic generator \(\phi\) of
Lemma~\ref{lem:flat-generic-generator}. At the embedding \(\sigma\)
write \(\Gamma,\Delta\) for the Gram matrices of the two frames.
The frame coefficient \(f\) then satisfies
\[
 |f|^2\,|\det\Gamma|/|\det\Delta|=c_\sigma^2
 \qquad(c_\sigma>0\text{ constant for the chosen }\phi).
\]
At a finite place \(w\) let \(\pi_w\) be a uniformizer. Degree zero on special-fibre curves
(Lemma~\ref{lem:flat-special-curve}) lets us apply
Lemma~\ref{lem:flat-vertical-ratio} to the projected determinant
lattice \(\det Q_\tau\). If \(R\) is the ring of a valued frame
over \(w\), the order computed in this lattice satisfies
\[
 \frac{\operatorname{ord}_{Q,R}(\phi)}{v_R(\pi_w)}
    =\max_{D\ {\rm vertical\ over}\ w}
       \frac{\operatorname{ord}_{Q,D}(\phi)}
            {\operatorname{ord}_D(\pi_w)} .
\]
To read the same section in the separate integral Lie frames,
Lemma~\ref{lem:projected-lie-frame-index} gives
\[
   \operatorname{ord}_{{\rm Lie},R}(\phi)
       =\operatorname{ord}_{Q,R}(\phi)-h_{R,\tau}.
\]
The Hodge multiplicities are fixed on this component, so
\(h_{R,\tau}/v_R(\pi_w)\) is the same at the CM frame and
at every other valued frame. The order and the metric may therefore
both be evaluated on the CM frame.
Lemma~\ref{lem:flat-cm-degree} evaluates their normalized frame
degree as \((p-q)\omega/2\).

For \(c\in F^\times\), replacing \(\phi\) by \(c\phi\) adds the
divisor of \(c\), whose normalized arithmetic degree is zero. Subtract
\(\delta_{w,\tau}\) from the order in each finite frame. Taking ideal
norms gives
\[
 \frac{1}{[F:\mathbb Q]}\sum_w\delta_{w,\tau}\log{\rm N}w
 =q\left(
 \frac{\log{\rm N}_{E/\mathbb Q}\mathfrak d_E}{[E:\mathbb Q]}
 -\log d_0\right)=q(R_E-D_0').
\]
Here the finite sum runs over the places of \(F\). The factor
\([F:E]\) from extending an ideal cancels in the normalization.
The corrected degree is the frame degree minus this sum. Its formula
contains no component-dependent quantity.
\label{cor:deep-cover-flat-degree-1}
\end{proof}

To compute the degree in Proposition~\ref{prop:deep-flat-line-degree}, we need one generator that can be evaluated in both kinds of frame. We first construct it on the generic component, then determine how its orders vary over a special fibre.

\begin{lemma}[Extension of the period generator]\label{lem:flat-generic-generator}
On the proper normal generic component of
Definition~\ref{hyp:flat-determinant-frames}, \(\mathcal L_\tau\) has an \(F\)-generator whose value in the prescribed frame at each complex embedding \(\sigma\) is a constant multiple of the normalized period determinant \(\Delta_{{\rm per},\sigma}\),
\[
 \phi_\sigma=c_\sigma\Delta_{{\rm per},\sigma},\qquad c_\sigma\ne0.
\]

\end{lemma}
\begin{proof}
In the marked period basis the first tangent block and the conjugate
of the second block have complementary directions. Pair the second
block by the Hermitian form and take the determinant. The marked
level comparison gives invariance under the change of period basis
on \(Z\). With \(\Gamma,\Delta\) the two frame Gram matrices, the
determinant identity, after division by the reference period constant,
reads
\[
 |f(v,w)|^2=\frac{|\det\Delta|}{|\det\Gamma|},\qquad
 \|f(v,w)e(v,w)\|^2=|f(v,w)|^2\frac{|\det\Gamma|}{|\det\Delta|}=1.
\]
It holds first at the embedding respecting the marked quadratic root.
Conjugating the frames preserves the absolute determinants.
\label{cor:deep-cover-flat-degree-2}

Take a proper smooth resolution of the generic component of the
normal model at \(2\), using the prime-to-\(2\) marking. In a
toroidal chart the boundary is \(z_1\cdots z_a=0\). The two
extending eigenframes are independent in the smooth semiabelian Lie
bundle. The boundary comparison for precisely these frames gives
\[
 C^{-1}Q(z)^{-A}\le u\le C Q(z)^A\quad (u=|\det\Gamma|,\ |\det\Delta|),
 \qquad Q(z)=1+\sum_{i=1}^a|\log|z_i|| .
\]
Consequently \(|f|,|1/f|\le C Q^A\), after increasing \(C,A\).
For \(g=f\) or \(1/f\), fix the other punctured coordinates \(z'\).
The Laurent coefficient of index \(-m\), \(m\ge1\), satisfies
\[
 \left|\frac{1}{2\pi i}\int_{|\zeta|=\rho}
          g(\zeta,z')\,\zeta^{m-1}\,d\zeta\right|
 \le C(z')\rho^m(1+|\log\rho|)^A\longrightarrow0 .
\]
Apply this to each boundary variable. The Laurent series on the
product of punctured discs has no negative multi-index in any
variable. Thus both \(f\) and \(1/f\) extend, and their product
remains one. Uniqueness on overlaps gives a generator on the resolution.
\label{cor:deep-cover-flat-degree-3}

Write \(Y\) for this proper geometrically integral resolution.
By GAGA and coherent base change,
\[
 H^0(Y,\mathcal L_\tau)\otimes_F\mathbb C
 \simeq H^0(Y_{\mathbb C}^{\rm an},\mathcal L_\tau^{\rm an})
 =\mathbb C\,\Phi ,
\]
where \(\Phi\) denotes the just-extended section and the line on
\(Y\) denotes its pullback. Any nonzero \(F\)-section therefore
generates after scalar extension, hence already over \(F\).
On the normal generic component use a local frame \(e\). Properness
of the resolution lifts each height-one valuation, so both
\(\phi/e\) and \(e/\phi\) have nonnegative order there. A normal
Noetherian domain is the intersection of its height-one valuation
rings. Thus both quotients are regular, as required. At another
finite prime the same construction on its prime-to-residue model
gives a generic generator. The quotient of the two generators lies
in \(F^\times\) since the generic component is proper and
geometrically integral. \label{cor:deep-cover-flat-degree-4}
\end{proof}

Fix the normal integral component of the prime-to-\(\ell\) model in
Definition~\ref{hyp:flat-determinant-frames}. Its generic point is
in the abelian locus. Let \(G\) be its smooth semiabelian group,
with its integral \(\mathcal O_E\)-action and rigidified polarizing
line. Write \(n=[E:\mathbb Q]\), \(d_G=\dim G\), and
\(r=p+q\), so \(2d_G=nr\). Over a local ring \(R\) on this
model choose formal coordinates at the identity and a rigidification
of the polarization biextension on the coordinate neighborhood.
These give the polarized tangent module
\[
       H_R=V_R\oplus V_R^*,\qquad V_R=\operatorname{Lie}(G)_R .
\]

The second summand is the space of linear terms of changes of
rigidification. Write \(U_a\) for the differential of \(a\in
\mathcal O_E\). The isomorphism between pulling the biextension
back by \(a\) in the first variable and by \(\bar a\) in the
second is represented in these frames by a unit equal to one on
both axes. Write \(C_a\) for its mixed linear coefficient, with
the first variable indexing the rows. The induced action is
\[
       (v,\alpha)\longmapsto
       (U_av,\ U_{\bar a}^{t}\alpha-C_a^{t}v).
\]

These formulas are integral even in characteristic two: the mixed
coefficient is the coefficient of \(st\) after substituting
\((sv,tw)\) with \(s^2=t^2=0\). Composition of the rigidified
isomorphisms gives matrix multiplication. A change of coordinates
acts on the first summand by its differential and on the second by
the inverse transpose, with a shear given by the mixed coefficient
of the change of frame. The inverse change has the same integral
description. Thus the displayed modules glue and commute with
base change.

\begin{auxlemma}[Local freeness over the integral coefficient algebra]
\label{lem:polarized-stalk-free}
On the marked component the polarized tangent module is locally
free of rank \(r\) over \(\mathcal O_X\otimes_{\mathbb Z}\mathcal O_E\).
In particular, on every geometric fibre of characteristic \(\ell\)
and for each Artinian local factor \(B\) of
\(k\otimes_{\mathbb Z}\mathcal O_E\),
\[
                   H_k\otimes_{k\otimes\mathcal O_E}B\simeq B^r .
\]
\end{auxlemma}
\begin{proof}
\proofstep{The abelian part of a geometric fibre.}
Work over an algebraically closed residue field \(k\) of
characteristic \(\ell\). In the interior the fibre is principally
polarized. At a boundary point choose a complete DVR dominating its stalk
whose generic point lies over the abelian generic point: normalize
the blowup of the maximal ideal, take the DVR of an exceptional
prime divisor, and complete. A one-dimensional normal stalk
itself supplies the DVR in dimension one. Its characteristic-zero
fraction field contains the local domain. After finite extension
the residue torus splits. The Raynaud extension and the cubical
line in the polarized degeneration datum reduce to
\[
  1\longrightarrow T\longrightarrow G_k
   \xrightarrow{\pi} A\longrightarrow1,\qquad
  L_{G_k}\simeq\pi^*L_A .
\]

Here the geometric fibre may be read after the further residue
field extension. Both ranks and local freeness descend faithfully.
For the period and character lattices \(Y\) and \(X^*(T)\), the
polarization induces an injection \(\Phi:Y\hookrightarrow X^*(T)\).
The degree formula for the degeneration datum gives
\[
   1=\deg\lambda_{\rm gen}
      =[X^*(T):\Phi Y]^2\deg\lambda_A,
      \qquad \deg\lambda_A=1 .
\]
The prime-to-residue model is used also when \(\ell\) divides the
original level. The Rosati adjoint is \(a\mapsto\bar a\).
\cite[Chapter~II, \S\S1--4, and Chapter~III, Corollary~7.2]{faltings-chai}

For \(0\ne a\in\mathcal O_E\) let \(K=\ker(a:A\to A)\).
The action is an isogeny when \(A\) has positive dimension, since
some integral multiple of \(a^{-1}\) is integral. The two
rigidified pullbacks of the Poincar\'e line give the perfect
Cartier pairing of \(K\) and
\(K^D=\ker(\bar a:A\to A)\).
\cite[Theorem~9.1]{milne-av-isogeny}
Let \(B_K\subset k[K]\) be the Hopf subalgebra generated by
the primitive functions, and set
\[
     J=\ker(K\longrightarrow\operatorname{Spec}B_K),\qquad
     t=\dim_k\operatorname{Lie}(K^D).
\]

A tangent vector to \(K^D\) is a character \(1+\epsilon u\) over
\(k[\epsilon]/\epsilon^2\). The character law reads
\[
             \Delta u=u\otimes1+1\otimes u .
\]
Thus the tangent vectors correspond to the primitive functions. The derivative of this function at a tangent
vector \(v\in\operatorname{Lie}K\) is the mixed coefficient of the
Cartier pairing. In our coordinates it equals
\(v^t C_a w\), up to the common choice of the inverse pairing, for
\(w\in\ker U_{\bar a}\). Indeed the two trivializations defining
the pairing on the kernels restrict to the rigidifications on the
axes. Evaluating their ratio on \((sv,tw)\) leaves just its
\(st\) coefficient.

A pair in the kernel on \(H_A\) therefore satisfies
\[
 U_av=0,\qquad U_{\bar a}^{t}\alpha=C_a^t v,\qquad
 v\in\operatorname{Lie}J .
\]
The last condition follows by pairing the middle equality with
each \(w\in\ker U_{\bar a}\). For a fixed \(v\) the ambiguity
in \(\alpha\) has dimension \(t\). For any finite group scheme
\(D\) over \(k\), its first Frobenius kernel has order
\(\ell^{\dim\operatorname{Lie}D}\), by the height-one
correspondence and the restricted enveloping algebra basis theorem.
\cite[Theorem~1.4.23]{brion-height-one}
Apply this first to \(J\), then to the Cartier dual of
\(\operatorname{Spec}B_K\), whose Lie algebra contains the \(t\)
independent primitives. Multiplicativity of orders in the finite
group scheme quotient gives
\[
   \dim\ker(a\mid H_A)\le \dim\operatorname{Lie}J+t
                          \le v_\ell(\deg a|_A).
\]
When \(A=0\) this contribution is zero.

\proofstep{The torus and the uniform kernel bound.}
The cocharacter lattice of \(T\) is torsion-free over
\(\mathcal O_E\), since multiplication by nonzero \(a\) divides
multiplication by a nonzero integer. It is projective over this
Dedekind ring. Put \(b_T=\operatorname{rk}_{\mathcal O_E}X_*(T)\)
and \(N_a=|N_{E/\mathbb Q}(a)|\).
The dimension of the kernel on \(\operatorname{Lie}T=X_*(T)\otimes k\)
is bounded by the sum of the \(\ell\)-adic elementary divisors
of multiplication by \(a\) on this lattice. For the toric term
and its conjugate dual the bounds are
\[
    \dim\ker U_{a,T}\le b_Tv_\ell(N_a),\qquad
    \dim\ker U_{\bar a,T}^t\le b_Tv_\ell(N_a) .
\]

Write \(P:V\to\operatorname{Lie}A\) for the quotient differential.
Use a rigidification pulled back from \(A\) on the geometric
fibre. The mixed coefficient there is \(P^t C_{a,A}P\).
In another rigidification the difference is a shear, so it
conjugates the action. The stable subspaces
\[
 \operatorname{Lie}T\oplus0\ \subset V\oplus P^*(\operatorname{Lie}A)^*\ \subset\ H_k
\]
have successive quotients \(\operatorname{Lie}T,H_A,
(\operatorname{Lie}T)^*\), the last with the conjugate action.
Kernel dimensions are subadditive along this filtration.
On the prime-to-\(\ell\) rational Tate module of \(A\), an
integral primitive element of \(E\) acts semisimply. Its rational
characteristic polynomial is a power of its minimal polynomial,
so every embedding occurs with multiplicity \(2\dim A/n\).
The degree of \(a\) is its determinant on this module.
\cite[Theorem~10.9 and Proposition~10.20]{milne-av-isogeny}
It follows that
\[
 \dim_k\ker(a\mid H_k)
   \le (2b_T+2\dim A/n)\,v_\ell(N_a)
   =r\,v_\ell(N_a).
\]

\proofstep{Recovering a free module at a ramified prime.}
On the characteristic-zero generic fibre the \(\tau\)-part of
\(H\) has dimension \(p+q\) for each embedding. At each stalk
the characteristic polynomial of \(a\) therefore equals that
of its regular action on \(\mathcal O_E\) raised to the \(r\)-th
power. The coefficients belong to the local domain and agree in
its fraction field, giving the integral identity
\[
      \det(X-a\mid H)=\det(X-a\mid\mathcal O_E)^{r}.
\]
For \(\mathfrak l\mid\ell\), of ramification
index \(e\) and residue degree \(f_{\mathfrak l}\), its factors
over \(k\) have the form \(B_\sigma=k[u]/(u^e)\). Each summand
\(H_\sigma\) has dimension \(er\) over \(k\). Indeed, choose
an element with nonzero residue generating the residue field at
\(\mathfrak l\) and zero residue at the other primes. The
multiplicities of its distinct nonzero eigenvalues in the preceding
characteristic polynomial are precisely these dimensions.

Choose \(a\) of valuation one at \(\mathfrak l\) and zero
at all other primes over \(\ell\), by the Chinese remainder
theorem. Then \(v_\ell(N_a)=f_{\mathfrak l}\). Multiplication
by \(a\) has the same kernel as \(u\) on the factors at
\(\mathfrak l\), and is invertible elsewhere. Thus
\[
 \sum_{\sigma}\dim_k\ker(u\mid H_\sigma)\le r f_{\mathfrak l},
 \qquad \dim_k\ker(u\mid H_\sigma)\ge \dim_k H_\sigma/e=r .
\]
All Jordan blocks have length at most \(e\), so equality forces
exactly \(r\) blocks of length \(e\) at every \(\sigma\).
This proves freeness over each \(B_\sigma\), including \(e>1\).
Over a characteristic-zero residue field the same conclusion follows
from the characteristic polynomial and separability of \(E\).

Freeness at the geometric fibre descends to projectivity at the
residue field by faithful flatness. The coefficient algebra there
is a product of local Artinian rings, and the common rank \(r\)
gives a free module. Lift a set of \(r\) coefficient generators
to the stalk. If \((b_i)\) is an integral basis of \(\mathcal O_E\),
the reductions of the \(b_i x_j\) are a basis of size \(nr=2d_G\).
Their determinant in the chosen tangent coordinates is a unit.
They give a basis over the local ring, hence the \(x_j\) are a
basis over its coefficient algebra.
\end{proof}

\begin{auxlemma}[Integral evaluation of a projected frame]
\label{lem:polarized-stalk-evaluation}
Let \(R\) be a stalk of this normal model and let the integral
character \(\sigma:\mathcal O_E\to R\) extend the label \(\tau\).
For any coefficient basis \(x_1,\ldots,x_r\) from
Lemma~\ref{lem:polarized-stalk-free}, the character quotient is
free with the indicated basis:
\[
 Q_{\sigma,R}=H_R/\sum_a(a-\sigma(a))H_R
             =\bigoplus_{j=1}^r R\,\bar x_j .
\]
Its generic-fibre realization is the image of \(H_R\) under
the characteristic-zero eigenprojection, and its bases and their
transition matrices specialize to the function field and every
valuation ring of every normalized special-fibre curve.
\end{auxlemma}
\begin{proof}
In coefficient coordinates, the map on each summand is
\(R\otimes\mathcal O_E\to R\), \(b\otimes a\mapsto b\sigma(a)\).
Its kernel is generated by \(1\otimes a-\sigma(a)\otimes1\).
After passing to the fraction field this evaluation is precisely the
corresponding projection in the split coefficient algebra. In
particular the integral quotient embeds there as the span of its
projected basis. No integral idempotent for the individual embedding
is required at a ramified prime.

Let \(C\) be a normalized curve, \(K=k(C)\) and
\(V_z=\mathcal O_{C,z}\). If \(x\) and \(x'\) are frames taken
at the image of \(z\) and of the generic point, compare them over
the latter stalk. Write their change of basis as \(T\) over its
coefficient algebra, in the convention \(x=x'T\). Then
\[
       \bar x=\bar x'\,T_\sigma,\qquad
       T_\sigma=(1\otimes\sigma)(T),\qquad
       \det T_\sigma\in\mathcal O_{X,\eta_C}^{\times}.
\]
The two maps from \(\mathcal O_{X,z}\) (at the image of \(z\))
to \(K\), through \(V_z\) and through \(\mathcal O_{X,\eta_C}\),
agree. Hence \(Q_{\sigma,V_z}\) in the generic curve frame
is \(T_\sigma V_z^r\), with the entries read in \(K\).
The determinant lattice is \((\det T_\sigma)V_z\). The determinant
can have nonzero \(z\)-valuation. Its valuation is exactly the
transition contribution to the degree on the curve.
\end{proof}

The integral image of an eigenprojection and the intersection with an
eigenspace may differ at a ramified prime. The following determinant
calculation compares the two conventions in
Proposition~\ref{prop:deep-flat-line-degree}.

\begin{auxlemma}[The local change from the projected to the two Lie frames]
\label{lem:projected-lie-frame-index}
Let \(R\) be the valuation ring of a finite extension of the local
coefficient field containing all embeddings of \(E\), write
\(K=\operatorname{Frac}R\), and let \(v_R\) be its integral valuation.
For a marked integral test write \(L=\operatorname{Lie}(G)_R\),
\(S_\gamma=L\cap(L_K)_\gamma\) and
\(n_\gamma=\dim_K(L_K)_\gamma\).
Let \(I_\tau\) be the embeddings inducing the same prime
\(\mathfrak l\mid\ell\) and residue embedding as \(\tau\).
For a uniformizer \(\theta\) of \(E_{\mathfrak l}\) put
\[
 h_{R,\tau}=\sum_{\gamma\in I_\tau\setminus\{\tau\}}
        n_\gamma v_R\bigl(\gamma(\theta)-\tau(\theta)\bigr).
\]
In the generic determinant identification, the lattice
\(\det S_\tau\otimes(\det S_{\bar\tau})^{-1}\) is contained in
\(\det Q_{\tau,R}\) with colength \(h_{R,\tau}\). In particular
the normalized difference of the two frame orders depends only on the
coefficient field, the place and the multiplicities.
\end{auxlemma}

\begin{proof}
On the factor selected by the common residue embedding the coefficient
algebra is monogenic. Indeed the unramified part splits over \(R\),
and a uniformizer generates the totally ramified integer ring over
its unramified part. On this factor put \(D=\theta-\tau(\theta)\).
It has the presentation
\[
 R[X]\big/\prod_{\gamma\in I_\tau}(X-\gamma(\theta)),\qquad
          Q_{\tau,R}=H_R/DH_R
\]
where \(H_R\) on the right denotes the selected factor. Take the
same factor of the exact sequence formed by the second summand and
the projection onto the first,
\[
         0\longrightarrow F\longrightarrow H_R
                              \longrightarrow L\longrightarrow0 .
\]
Here \(L\) in the sequence also means its selected factor, and
\(F\) is the corresponding factor of the integral dual of the
full Lie module, with conjugate action. Its projected image is
\(S_{\bar\tau}^*\): restriction to the saturated free sublattice
\(S_{\bar\tau}\) is surjective over a DVR, and the generic
restriction is the dual eigenprojection.

Since \(Q_{\tau,R}\) is free, the image of \(F/DF\) in it is its
torsion-free quotient, which the preceding description identifies with
\(S_{\bar\tau}^*\). Right exactness gives
\[
     0\longrightarrow S_{\bar\tau}^*\longrightarrow Q_{\tau,R}
                       \longrightarrow L/DL\longrightarrow0 .
\]
Write \(I=L\cap\bigoplus_{\gamma\ne\tau}(L_K)_\gamma\) on this
factor. The torsion in \(L/DL\) is \(I/DL\) and the free quotient
is \(L/I\). The colength in the assertion is consequently
\[
    \operatorname{length}(I/DL)+
    \operatorname{length}\bigl(L/(I+S_\tau)\bigr).
\]

This follows by a basis change over the DVR, first
saturating \(S_{\bar\tau}^*\) in \(Q_{\tau,R}\) and then taking
the index of \(S_\tau\) in the free quotient \(L/I\).
The map \(D\) identifies \(L/(I+S_\tau)\) with \(DL/DI\),
since its kernel on \(L\) is \(S_\tau\). The two lengths add to
\[
  \operatorname{length}(I/DI)
     =v_R\bigl(\det(D\mid I_K)\bigr)=h_{R,\tau}.
\]
The formula holds also for an unramified factor, when the sum is empty.
Passing to a further finite extension multiplies both the colength
and \(v_R(\pi_w)\) by its ramification index.
\end{proof}

\begin{lemma}[Degree on special-fibre curves]\label{lem:flat-special-curve}
Use the integral eigenframes of Definition~\ref{hyp:flat-determinant-frames}
on the normal model at a finite place of residue characteristic \(\ell\).
The determinant line measuring the order of the section \(\phi\)
of Lemma~\ref{lem:flat-generic-generator} has degree zero on the normalization \(\widetilde C\) of every complete integral curve \(C\) in this special fibre,
\[
 \deg_{\widetilde C}(\mathcal L_\tau|_{\widetilde C})=0.
\]

\end{lemma}
\begin{proof}
Apply Lemma~\ref{lem:polarized-stalk-evaluation} on the marked
model after extending constants to contain the integral label
\(\sigma\). On the generic fibre the projected polarized tangent
space is an extension of \(V_\tau\) by \(V_{\bar\tau}^*\).
Lift the \(v_i\) to this extension and adjoin the dual frame of
the \(w_j\). Its determinant is independent of the lifts and
transforms by \(\det A/\det B\). It identifies
\[
       \det Q_{\sigma,\operatorname{Frac}R}
                 =\det V_\tau\otimes(\det V_{\bar\tau})^{-1} .
\]
Its extension on the curve has the integral determinant lattices
computed in Lemma~\ref{lem:polarized-stalk-evaluation}. These
are the orders to which Lemma~\ref{lem:flat-vertical-ratio} will
be applied. The comparison with the separate Lie frames in the CM
calculation keeps the different at the given place.
\label{cor:deep-cover-flat-degree-5}

\proofstep{Compare the two Hodge degrees on a Siegel model.}
Enlarge the finite field of constants \(k\) to split the residue
fields of \(\mathcal O_E\) above \(\ell\). Write \(e\) for a primitive
idempotent in \(k\otimes\mathcal O_E\), and write \(H_e=eH\).
We compare determinants using the underlying principally polarized
family. Choose a full Siegel marking of some order \(n\ge3\) prime
to \(\ell\) on a finite cover of the abelian generic locus. Take the
normalized closure of the graph of its map to a projective Siegel
toroidal compactification at that level. If \(X\) denotes our
model and \(X'\) this graph, the projections give
\[
 X\ \longleftarrow\ X'\ \longrightarrow\ \overline{\mathcal A}_{d_G,n}^{\rm tor}.
\]
The first projection is proper and surjective. On this normal cover the two
semiabelian families agree: their generic isomorphism and its inverse
extend uniquely. \cite[I, Proposition~2.7]{faltings-chai}

There is a complete integral curve upstairs dominating
\(\widetilde C\) after normalization. Indeed a closed point of the
preimage of its generic point has residue field finite over
\(k(\widetilde C)\); take the closure in the proper fibre product
and normalize. For this finite map \(h:C'\to\widetilde C\),
\[
              \deg_{C'} h^*\det H_e
                    =(\deg h)\deg_{\widetilde C}\det H_e ,
\]
using degrees relative to the same ground field. We can therefore
compute on \(C'\), even for a curve contained in the boundary.

Let \(D\) be the pullback of the covariant logarithmic de Rham
bundle on the Siegel compactification. Its Hodge and conjugate
filtrations form an \(F\)-zip: the two filtrations are by subbundles
and their graded bundles match in reverse order after Frobenius
pullback. The boundary assertion here is the logarithmic one for
principally polarized Siegel families in characteristic \(\ell\),
including \(\ell=2\). In a toroidal chart its covariant realization
is the Lie algebra of the period-equivariant universal vector
extension of the Raynaud extension. For split torus \(T\), abelian
quotient \(A\) and period lattice \(Y\), its weight filtration has
successive quotients
\[
          \operatorname{Lie}T,\qquad D(A),\qquad Y\otimes\mathcal O .
\]
Here \(D(A)\) denotes the ordinary de Rham homology of \(A\). \cite[VI, Theorem~1.1]{faltings-chai}
\cite[\S2D, pp.~476--477]{vdg-looijenga-boundary-zip}

The coefficient action on these filtrations also extends. On the
completed normal degeneration chart, after splitting the torus,
the action preserving the polarization in the Rosati sense acts
integrally on the Raynaud extension and its period homomorphism
\(u\). If \(a_Y,a_{G^\natural}\) are the induced maps, then
\[
 u:Y\longrightarrow G^\natural_\eta,\qquad
          u\circ a_Y=a_{G^\natural}\circ u .
\]
This is the equivalence for integral endomorphism structures in
\cite[Theorem~5.1.1.4]{lan-pel-degeneration}.
Applying the universal vector extension and its Lie functor gives
the action on \(D\). It commutes with relative Frobenius and
Verschiebung by functoriality on the degeneration data. The
actions agree on overlaps by their common generic action. In
particular the construction applies to every \(a\in\mathcal O_E\),
including elements of positive \(\ell\)-adic valuation.

Put \(V=\operatorname{Lie}(G)|_{C'}\), and let \(U=V^*\) with the
conjugate dual action induced by principal polarization. Write
\(F\) for absolute Frobenius of \(C'\) and \(F_k(e)\) for the
idempotent obtained by raising the scalar coefficients of \(e\)
to the \(\ell\)-th power. Taking the indicated direct summands
of the two filtrations gives exact sequences on the whole curve:
\[
\begin{aligned}
 0&\longrightarrow U_e\longrightarrow D_e
                 \longrightarrow V_e\longrightarrow0,\\
 0&\longrightarrow F^*V_e\longrightarrow D_{F_k(e)}
                 \longrightarrow F^*U_e\longrightarrow0 .
\end{aligned}
\]
The quotients are vector bundles, so pullback to the curve and
projection by the idempotents preserve exactness.

The polarized
tangent bundle itself has \(0\to U_e\to H_e\to V_e\to0\).
Consequently degree additivity and \(\deg F^*L=\ell\deg L\)
for a line bundle give
\[
 \deg H_{F_k(e)}=\deg D_{F_k(e)}
       =\ell(\deg U_e+\deg V_e)=\ell\deg H_e .
\]
The permutation \(e\mapsto F_k(e)\) has finite order \(s>0\).
Iteration gives \((\ell^s-1)\deg H_e=0\). This argument uses
the whole local Artinian factor selected by \(e\), before reduction
by its maximal ideal.

The local Artinian factor \(B\) at the chosen character has residue
field \(k\). Filter \(B\) by ideals with simple quotients \(k\).
If \(b=\operatorname{length}_B B\), Lemma~\ref{lem:polarized-stalk-free} identifies
each quotient of the induced filtration of \(H_e\) with the same
bundle \(H_e\otimes_B k\) (up to a constant one-dimensional factor).
By degree additivity,
\[
 0=\deg H_e=b\,\deg(H_e\otimes_B k).
\]
The determinant of the eigencharacter quotient therefore has degree
zero, including when \(B\) has nilpotents.
\label{cor:deep-cover-flat-degree-6}
\end{proof}

\begin{auxlemma}[A curve through a prescribed point]
\label{lem:flat-specialization-curve}
Let \(T\) be an integral proper variety of positive dimension over
a finite field, \(U\subset T\) a dense open subset, and \(z\) a closed
point of \(T\). Let \(k\) be the finite ground field. There is a finite extension \(k'/k\), a complete integral curve \(C\subset T_{k'}\) and a point \(z'\in C\) over \(z\) such that its generic point satisfies
\[
 z'\in C\cap(T_{k'})_z,\qquad \eta_C\in U_{k'}.
\]

\end{auxlemma}
\begin{proof}
Over an algebraic closure take a reduced irreducible component through
a lift \(z'\) of \(z\). The inverse image of \(U\) is dense there.
Take an affine neighborhood \(V\) of \(z'\) on this component and a
finite Noether normalization \(f:V\to\mathbb A^d\).
For the closed subset \(Z=V\setminus U_{\bar k}\) one has
\[
   \dim f(Z)=\dim Z<d\qquad(Z\ne\varnothing).
\]
Choose a line through \(f(z')\) and a point outside \(f(Z)\).
For any component \(C_0\) of its inverse image through \(z'\),
the \(d-1\) linear equations and finiteness give
\[
 1=d-(d-1)\le \dim C_0=\dim f(C_0)\le1 .
\]
The finite image is the whole line, so \(C_0\) meets the given open
set. Its closure in the proper variety is the required integral curve.
The component, line and curve are defined by finitely many data over
the algebraic closure of a finite field and hence descend to a finite
extension.
\end{proof}

\begin{lemma}[The vertical order on a connected fibre]\label{lem:flat-vertical-ratio}
Let \(\mathcal L\) be the integral determinant line on the proper
normal component over a DVR with uniformizer \(\pi\) and connected
special fibre. Assume it has degree zero on every normalized
complete integral curve in that fibre. If \(\phi\) generates on
the generic fibre, then, for every valued frame over this DVR,
\[
 \operatorname{ord}(\phi)=\frac{a}{m}\operatorname{ord}(\pi),\qquad
 \frac{a}{m}=\max_{D\ {\rm vertical}}
            \frac{\operatorname{ord}_D(\phi)}{\operatorname{ord}_D(\pi)},
 \quad m>0 .
\]
\end{lemma}
\begin{proof}
The rational section \(s=\pi^a\phi^{-m}\) of \(\mathcal L^{-m}\)
has horizontal order zero. Its order along a vertical prime \(D\) is
\[
 a\,\operatorname{ord}_D(\pi)-m\,\operatorname{ord}_D(\phi)\ge0 .
\]
By normality \(s\) is regular. Its unit locus \(U\) contains the
generic point of a component attaining the maximum.

Take any irreducible special-fibre component \(T\) meeting \(U\).
If \(T\setminus U\) is nonempty, choose a closed point \(z\) there.
If \(\dim T=0\) then \(T\) consists of a single point, contrary to
the choice of \(z\). Otherwise apply Lemma~\ref{lem:flat-specialization-curve}
to \(T_{\rm red}\) and its open unit locus, and normalize the resulting
curve. Its image in the original fibre is an integral curve. Pullback
along the finite map of normalized curves preserves degree zero.

On this complete normal curve \(C\), the pullback of \(s\) is
regular and generically nonzero, so
\[
 0=-m\deg_C\mathcal L=\deg_C\operatorname{div}(s)
       =\sum_{x\in C_{\rm cl}}[k(x):k_C]\operatorname{ord}_x(s),
 \qquad \operatorname{ord}_x(s)\ge0 ,
\]
where \(k_C\) is its constant field. A point over \(z\) would
contribute a positive summand. Hence \(T\subset U\). Across any
intersection with a neighboring component its generic point belongs
to the same open \(U\). Connectedness of the finite component graph
gives the result on the whole fibre. Pulling the unit \(s\) back to
a valued frame gives \(m\operatorname{ord}(\phi)=a\operatorname{ord}(\pi)\).
Changing the generic generator by \(c\in F^\times\) translates its
order by \(\operatorname{ord}(c)\). The local models and framed
uniformization supply nonempty tests at the required places.
\label{cor:deep-cover-flat-degree-7}
\end{proof}

The vertical calculation has made the normalized orders independent of the component of a connected fibre. We can therefore evaluate the remaining degree at the prescribed CM frame.

\begin{lemma}[CM calculation in the frame metrics]\label{lem:flat-cm-degree}
For the polarized determinant line on \(Z\) in
Definition~\ref{hyp:flat-determinant-frames}, before the finite-order
correction of Proposition~\ref{prop:deep-flat-line-degree}, the normalized arithmetic degree in the CM frame norms satisfies
\[
 \widehat{\deg}_{\rm CM}(\mathcal L_\tau)=\frac{p-q}{2}\,\omega.
\]

\end{lemma}
\begin{proof}
At a pure CM test on the component use the two elliptic types of
\(\mathbb Q(\sqrt{-d_0})\), with their good integral abelian models,
integral Lie generators and normalized principal polarizations. Write
\(n_{\pm,s}\) for the squared Lie lengths at a complex embedding \(s\).
In a common finite field of definition \(k/F\), write \(\langle
\cdot\rangle_\infty\) for the mean over all \([k:\mathbb Q]\) complex
embeddings, counting conjugates. The Chowla--Selberg comparison and
the functional equation of the odd primitive quadratic character give
\[
 \begin{aligned}
 -\langle\log n_{\pm,s}\rangle_\infty=X
 &=\gamma+\log(2\pi/\sqrt{d_0})-L'(1,\chi)/L(1,\chi)\\
 &=L'(0,\chi)/L(0,\chi)+D_0'/2,\qquad
 \omega=X-R_E+D_0'/2 .
 \end{aligned}
\]
The mean is invariant under further field extension.
\cite{colmez-chowla-selberg}\label{cor:deep-cover-flat-degree-8}

Let \(T\) change the reference period basis to the two orthogonal
CM blocks. Let \(G\) be the reference Hermitian determinant, and
\(G_+,G_-\) the two block determinants. In these frames the
coefficient of \(\phi\) is \(\eta\,\tau(G_-)/\tau(\det T)\).
The elliptic period scalars for the two blocks are included in
\(\eta\). Evaluating the algebraic generator on the algebraic
frames shows that \(\eta\in k^\times\) after finite extension.
The identity \(G_+G_-=(\det T)\overline{\det T}\,G\) cancels the
transition determinant against the determinant metric. With the
two squared length scales \(n_{\pm,s}\sqrt{d_0}\), this gives
\[
 A_s=\log|s\eta|+\tfrac12\log|s(\tau G)|
       +\tfrac p2(\log n_{+,s}+D_0'/2)
       -\tfrac q2(\log n_{-,s}+D_0'/2).
\]
\label{cor:deep-cover-flat-degree-9}

At \(\ell\) take a locally integral CM basis with
\(t_\ell=\det T_\ell\). For \(j:k\hookrightarrow\mathbb C_\ell\)
put \(i=j\circ\tau\), \(D_j=\log|i(\mathfrak d_E)|_\ell\) and
\(d_\ell=\log|d_0|_\ell\). The absolute value of an ideal here is
that of a local generator. Unimodularity of the polarized period
lattice and trace duality with the different give
\[
 \log|i(t_\ell)|_\ell+\log|i(\bar t_\ell)|_\ell+\log|i(G)|_\ell
       =-(p+q)D_j,\qquad \log|i(G_-)|_\ell=-qD_j .
\]

The second identity uses the orthogonal pure-type block. The
locally integral period basis spans the CM product lattice over
the marked quadratic order up to prime-to-\(\ell\) index.
Projecting to each integral eigenlattice by the trace-dual basis
contributes a generator of \(\mathfrak d_E/(\sqrt{-d_0})\), whose
log absolute value is \(D_j-d_\ell/2\) on either conjugate
label. The two frame exponents are \(-p,+q\), giving the finite log norm
\[
 B_j=\log|j\eta|_\ell-\log|i(t_\ell)|_\ell-qD_j
                              +(q-p)(D_j-d_\ell/2).
\]
\label{cor:deep-cover-flat-degree-10}

To sum, put \(B'_j=B_j-q(D_j-d_\ell)\). Conjugation permutes
the embeddings of \(E\) with equal multiplicities at \(\ell\),
so the two sums for \(t_\ell,\bar t_\ell\) coincide. The first
trace-dual identity yields
\[
 \sum_{j\mid\ell} B'_j
 =\sum_{j\mid\ell}\left(\log|j\eta|_\ell+\tfrac12\log|i(G)|_\ell
              -\tfrac{p+q}{2}D_j+\tfrac{p+q}{2}d_\ell\right).
\]
Write \(\langle\cdot\rangle_f\) for the sum over finite embeddings,
divided by \([k:\mathbb Q]\), with \(|\ell|_\ell=\ell^{-1}\).
Only finitely many places contribute. Ideal norms give
\(\langle D_j\rangle_f=-R_E\), \(\langle d_\ell\rangle_f=-D_0'\).
The product formula cancels the \(\eta,G\) terms. Therefore
\[
 \begin{aligned}
 -\langle A_s\rangle_\infty-\langle B'_j\rangle_f
 &=\tfrac{p-q}{2}(X-D_0'/2)+\tfrac{p+q}{2}(D_0'-R_E)\\
 &=\tfrac{p-q}{2}\omega+q(D_0'-R_E),\\
 -\langle A_s\rangle_\infty-\langle B_j\rangle_f
 &=\tfrac{p-q}{2}\omega .
 \end{aligned}
\]
The last equality restores \(B_j-B'_j=q(D_j-d_\ell)\) in the
frame norms. \label{cor:deep-cover-flat-degree-11}
\end{proof}


\begin{lemma}[Slope change under the determinant twist]
For the flat twist from \(\Lambda_M\) to \(\Lambda_-\),
\[
 \mu(\Lambda_-)-\mu(\Lambda_M)\ge-M_i-(11C_A+83)U .
\]
\end{lemma}
\begin{proof}
The degree calculation gives
\[
 c_{1,\tau}=-\tfrac12(p-q)\omega
           -q\bigl(\log d_0-\log\operatorname{rd}(E)\bigr).
\]
Since \([E:\Q]=2f\), the discriminant bound reads
\[
       0\le \log\operatorname{rd}(E)
       =\frac{\log|\operatorname{disc}E|}{2f}\le2L_{\rm sc}.
\]
Since \(0\le\log d_0\le L_{\rm sc}\), the last term in the
identity has absolute value at most \(2qL_{\rm sc}\).
Together with \(|\omega|\le3L_{\rm sc}\) and \(p+q=r\), this gives
\[
 |c_{1,\tau}|\le(3|p-q|/2+2q)L_{\rm sc}
          \le(7/2)rL_{\rm sc}\le9rL_{\rm sc}.
\]
\label{cor:covolume-component-transfer-10}

In the bad-prime allowance every prime divides
\(15g_0g_1 C_E\), where \(C_E\) is the conductor
parameter containing the prime divisors of the CM field discriminant, including \(2\) and \(d_0\).
This uses \(N_{\rm lev}=g_0g_1\,15^e\) and auxiliary level one.
Thus the sum of the logarithms of the distinct bad primes
is at most \(\log(15g_0g_1 C_E)\). The parameter bound
\(C_E^{36}t^{150}\le b^{1/2}\) with \(t\ge1\) and
\(\log b\le8L_{\rm sc}\) implies \(\log C_E\le L_{\rm sc}\).
The two bounds \(\log g_k\le9L_{\rm sc}+1\) now give
\[
 \sum_{p\ {\rm bad}}\log p
 \le\log 15+\log g_0+\log g_1+\log C_E
 \le 19 L_{\rm sc}+2+\log 15\le20L_{\rm sc},
\]
since \(L_{\rm sc}\ge10\). This proves the allowance \(40L_{\rm sc}\)
used below on the chosen deep level even as \(e\) grows.
\label{cor:covolume-component-transfer-11}

Combining the flat-degree estimate with the bounds for the shifts \(w_\tau\) gives
\[
 \begin{split}
 |M_i-(p-q)\omega\textstyle\sum_\tau w_\tau|&\le C_A U,\\
 q L_{\rm sc}\textstyle\sum_\tau|w_\tau|&\le C_A U,\qquad
 |c_{1,\tau}+(p-q)\omega/2|\le3qL_{\rm sc}.
 \end{split}
\]
The even column count is \(2m_0\). The combined weight of
the base line, the parallel shift by \((w_\tau)_\tau\) and the
\(m_0\)-th theta root has exactly the exponents of the
negative-shift line. Indeed this parallel weight is subtracted from the
base weight in the definition of the unshifted quotient and
is restored in the parallel twist. The finite valuation rules,
cusp conditions, metric factor at every embedding and hence
the Petersson norms depend on these combined exponents.

The flat-line twist preserves rank. Its normalized degree
increment per rank in the squared metric is bounded below by
\[
 2\left(\sum_\tau w_\tau c_{1,\tau}
       -q\ell_d\sum_\tau|w_\tau|
       -m_0fr\{\tfrac12\log(2\pi)+\ell_d\}-\ell_S\right)
\]
on the unramified deep charts. Here \(\ell_d\le2L_{\rm sc}\)
is the different allowance and \(\ell_S\le40L_{\rm sc}\)
the sum of bad-prime logarithms. Twice the first sum is at
least \(-M_i-7C_AU\) by the three preceding inequalities for \(w_\tau\) and \(c_{1,\tau}\).
The second loss costs at most \(4C_A U\). Using
\(2m_0\le n_V\), \(frn_V L_{\rm sc}\le b\le U\) and
\(10\le L_{\rm sc}\le U\), the last two losses cost at most
\(3U+80U\). Adding the four losses gives the stated lower bound.
Write \(C_F=11C_A+83\). \label{cor:covolume-component-transfer-12}
\end{proof}

\begin{auxlemma}[Simultaneous evaluation above an unramified point]
\label{lem:adjoint-etale-fibre}
Let \(f:M\to Y\) be a generically finite surjective morphism of smooth
projective complex varieties of common dimension \(d\ge1\). Let \(A\)
be a nef and big integral line bundle on \(Y\), and \(P\) a nef line
bundle on \(M\). Choose \(y\) in the finite \'etale locus such that
\(A^{\dim V}\cdot V>0\) for every positive-dimensional integral
\(V\ni y\). Put \(h=d(d+1)/2\). For \(m\ge h+1\), the evaluation
on the entire reduced fibre is surjective:
\[
 H^0(M,L)\longrightarrow\bigoplus_{z\in f^{-1}(y)} L(z),
 \qquad L=K_M+m f^*A+P .
\]
\end{auxlemma}
\begin{proof}
Each positive intersection at \(y\) is an integer and is at least
one. On \(Y\) use the rational divisor \(N=(h+1/2)A\). For a
subvariety of dimension \(i\) through \(y\) one has
\[
       N^i\cdot V\ge(h+1/2)^i>h^i,\qquad
       \sum_{i=1}^d i/h=1 .
\]
Koll\'ar's isolated non-klt point theorem, applied to the smooth
pair \((Y,0)\), gives an effective rational divisor \(D\equiv N\)
whose non-klt locus near \(y\) consists of \(y\) alone.
\cite[Theorem~6.4]{kollar-pairs-adjoint}

Put \(J=\mathcal J(M,f^*D)\), the multiplier ideal. Discrepancies
and multiplier ideals pull back under an \'etale morphism, as one
sees by pulling back a log resolution over the \'etale open.
Consequently the points of \(f^{-1}(y)\) are isolated components of
the support of \(\mathcal O_M/J\). Since \(f\) is generically finite,
\(f^*A\) is nef and big. Moreover
\[
    L-K_M-f^*D\equiv (m-h-1/2) f^*A+P
\]
is nef and big. Nadel vanishing gives \(H^1(M,L\otimes J)=0\).
Thus every section of \(L\otimes(\mathcal O_M/J)\) lifts to \(M\).
The Artinian local scheme at each isolated point is open and closed
in this support. A chosen residue value lifts there, and one may
take zero on all the other components. This gives all the indicated
values simultaneously. \cite[Theorem~2.16 and proof of Theorem~5.8]{kollar-pairs-adjoint}
\end{proof}

The twist has put the slope at the desired weight. Descent to the selected level next
requires a section whose translates can be separated by character traces. Let \(S\) be
the finite level group, acting on the right, and let \(\chi\) be the prescribed
character.
\begin{lemma}[A short section separating the translates]
\label{lem:deep-short-orbit-section}
For the selected Hodge shift \(W\) there is a short integral
section \(v\) such that \(\det(v^{hg})_{h,g\in S}\) is a
non-zero-divisor. Its log supremum norm satisfies
\[
 \log\|v\|_{\sup}=O\big((frL_{\rm sc})^9\big)
\]
provided the component generating families have size \(N_c\),
Petersson bound \(\Lambda\) and pointwise factor \(\beta\) with
\(\log N_c,\log\beta=O(z_0^3)\), \(\log\Lambda=O(z_0^9)\) and
\(z_0=frL_{\rm sc}\), at the prescribed deep level of
Lemma~\ref{hyp:covolume-deep-component}.
\end{lemma}
\begin{proof}
The same section must give a nonzero determinant on every component.
The chosen Hodge exponent \(c_0\) exceeds
\(r+\binom{d_X+1}{2}\), with \(d_X\) the dimension of the
moduli component. At a complex framed point over its generic
point consider its \(S\)-orbit. The marking at a prime
greater than \(r+1\) is selected as follows. The deep level
\(N_{\rm lev}\) is divisible by \(g_0g_1\). In the marked type frame the first
level subgroup is upper unitriangular at both type primes.

For the other subgroup, indexed by \(k\), use the type prime
with index different from \(k\), where it is upper
unitriangular. At a place over the chosen \(g=g_j\), the
half-level marking determines an ideal \(\mathfrak I\) containing
the marking ideal. The complete splitting and squarefree
support of \((g)\) give
\[
  \mathcal O_E/\mathfrak I=\mathbf F_g,\qquad
  g\notin\mathfrak I^2,\qquad g>r+1.
\]
In the marked reduction the matrices are unitriangular. Hence
\[
             (s-I)^r\equiv0\pmod{\mathfrak I}
                   \qquad (s\in S).
\]

If two translates have the same underlying complex moduli
point, their relative level change fixes that point. Pulling
the universal group and the level isomorphism back to it
gives a finite-order \(\mathcal O_E\)-linear automorphism
\(\alpha\) of its semiabelian fibre. At \(\mathfrak I\) the
Tate-module identification reads its reduction on
\(\mathfrak I\)-torsion by the marked action, up to inversion.
Here the \(\mathfrak I\)-torsion is the part of the marking
obtained by multiplying by a generator of
\(\mathfrak m_b\mathfrak I^{-1}\) at \(\mathfrak I\),
where \(\mathfrak m_b\) is the marking ideal. Its dimension
is at most \(r\); the corresponding free
\(\mathcal O_{E,\mathfrak I}^{\wedge}=\mathbf Z_g\) Tate module
has rank at most \(r\). Every eigenvalue \(\zeta\) of \(\alpha\)
is torsion and reduces to one. Its order is therefore a power
of \(g\). A nontrivial one would satisfy
\[
 g-1\le[\mathbf Q_g(\zeta):\mathbf Q_g]
       \le\operatorname{rk}_{\mathbf Z_g}T_{\mathfrak I}\le r<g-1 .
\]
Thus all eigenvalues equal one; finite order gives semisimplicity
in characteristic zero.

To see faithfulness on the fibre, use the period
lattice of its complex exponential map. Its rational span
is a vector space over \(E\). Tensoring with \(E_{\mathfrak I}\)
gives the rational Tate module just used, so the action on
the period lattice itself is trivial. The period lattice of
a complex semiabelian variety spans its Lie algebra over
\(\C\), by the torus and abelian quotient. Thus \(\alpha\) acts as the identity on the fibre and hence on its
level lattice quotient. The marking is faithful. If \(x\) is the
framed point, this proves the required orbit separation:
\[
                x^g=x^h\quad\Longrightarrow\quad g=h
                              \qquad(g,h\in S),
\]
where equality means equality of the underlying moduli points.

Let \(X\) be the smooth projective toroidal moduli component, of
dimension \(d_X\), and \(\Gamma\subseteq S\) its stabilizer. Taking
the closure of the graph of all its translates and normalizing gives
an equivariant modification \(X'\to X\) and a finite quotient
\(q:X'\to X'/\Gamma\). These modifications are isomorphisms over
the common interior. Writing \(Z^\circ\) for the interior of this component gives
\[
 X'|_{Z^\circ}=Z^\circ,\qquad
 (X'/\Gamma)|_{Z^\circ/\Gamma}=Z^\circ/\Gamma .
\]

The Hodge line \(\lambda\) has the linearization induced by
pullback of invariant differentials. The level isomorphisms on the
dense interior extend to the semiabelian models and satisfy the group
law by uniqueness. \cite[I, Proposition~2.7]{faltings-chai}
A stabilizer at a boundary point acts on its semiabelian fibre.
The same marked Tate-module argument above makes this action trivial,
hence also its action on the determinant of invariant differentials.
Thus \(\lambda\) descends at the first power.

Indeed, near an orbit choose a local generator at one point. Prescribe
its translates at the other points and lift by the Chinese remainder
theorem on an affine inverse image of the quotient. Averaging the
lift over \(\Gamma\) gives an invariant section with the same values,
because the point stabilizers act trivially on the fibres. It is a
generator on an invariant neighborhood of that orbit. The invariant
transition units descend and define a line \(A\) with
\[
                 q^*A=\lambda|_{X'} .
\]

The notation on the right denotes pullback from \(X\).
The full Hodge determinant \(\omega\) comes from the ample line
on the minimal compactification, whose contraction is an isomorphism
on the open moduli space. On the full flag its weight is
\((1,\ldots,1;-1,\ldots,-1)\) at each CM label. Our Hodge
shift has weight \((0,\ldots,0;-1,\ldots,-1)\). Thus
\[
  2\lambda-\omega \quad\text{has constant weight }-1
            \quad\text{at each label}.
\]
Such a weight and its negative are dominant. The two curve
inequalities identify its numerical class with zero. Since
\(\omega\) is big and nef, the same holds for \(\lambda\).
\cite[Remark~16.2(2)]{lan-ramified-models}

Resolve the quotient and the graph over it, with both modifications
isomorphisms over a common dense open. Denote the induced morphisms
by \(f:M\to Y\) and \(a:M\to X\), writing \(A\) also for the
pullback to \(Y\). Then
\[
 f^*A=a^*\lambda,\qquad
 A^i\cdot V>0\quad(\dim V=i>0,\ y\in V)
\]
at a general quotient point \(y\). Indeed nefness descends by the
projection formula on curves. Writing \(D=\deg f\) gives
\[
         D\,A^{d_X}=(f^*A)^{d_X}=\lambda^{d_X}>0 .
\]
Write \(a_0 A=H+E\) with \(H\) ample and \(E\) effective.
For \(y\notin\operatorname{Supp} E\), the mixed intersections on
each \(V\ni y\) satisfy
\[
     (a_0 A)^i\cdot V\ge H^i\cdot V>0 .
\]
In characteristic zero the generically finite map is finite \'etale
after shrinking this open. The framed point was chosen above the
generic point over the number field, so it and its translates belong
to these opens after the chosen complex embedding.


Write \(D_\partial\) for the toroidal boundary and
\(\kappa=K_X+D_\partial\) for its logarithmic canonical line. With
\(m=1+\binom{d_X+1}{2}\), the chosen \(c_0\) has \(c_0-m\ge r\).
The pullback of \(P_X=(c_0-m)\lambda-\kappa\) to the full flag
has dominant weight: subtracting the logarithmic canonical weight
uses at most \(r\) of the Hodge gap. The curve criterion of
Section~\ref{subsec:component-rank-comparison} and surjectivity of
the flag projection therefore give \(P_X\) nef. We have
\[
 K_X+m\lambda+P_X=c_0\lambda-D_\partial.
\]

Apply Lemma~\ref{lem:adjoint-etale-fibre} on \(M\) with \(P=a^*P_X\).
The regular top forms on a smooth projective modification push
forward to the canonical bundle on \(X\): they agree off the
exceptional images, of codimension at least two, and extend there
in the locally free canonical sheaf. Pullback of regular top forms
gives the reverse inclusion. Hence its sections descend at this
same weight and their values at the orbit points agree up to
invertible frame changes.

The canonical boundary section converts these section values into values of cusp forms.
Indeed, if \(P_1,\ldots,P_t\) are the generic orbit points in
this component and \(b_\partial\) is the boundary section, then in
the chosen frames
\[
 \operatorname{ev}(b_\partial s)
   =\operatorname{diag}\bigl(b_\partial(P_1),\ldots,b_\partial(P_t)\bigr)
                        \operatorname{ev}(s).
\]
Each \(P_i\) lies off the boundary, so this diagonal map is an automorphism of the space of section values. These are the
geometric separation and weight hypotheses used here.
\label{cor:covolume-component-transfer-13}


Extend the component sections by zero. At the framed test point
\(x_Z\) chosen over each component \(Z\), their values span the
delta function at the identity of \(S\). For a section \(v\)
with this evaluation vector the group determinant at \(x_Z\) is
\[
 \det\big(v^{hg}(x_Z)\big)_{h,g\in S}
                =\det\big(\mathbf1_{hg=1}\big)_{h,g\in S}=\pm1 .
\]
Thus the determinant section is nonzero at the generic point
of \(Z\); write \(D_Z\) for its polynomial in the coefficients
of \(v\) at that point.

Let \(F\) be the common field of definition of the marked level and
its section space. Extension of scalars in the equalizer of sections
on a finite affine cover gives
\[
 H^0(X_F,\mathcal L)\otimes_{F,\iota}\mathbb C
     \simeq H^0(X_{\mathbb C},\mathcal L_{\mathbb C}).
\]
Thus a finite rational section space supports these tests for all
\(Z\). The product \(\prod_Z D_Z\) on it is nonzero, and
rational tuples over the infinite coefficient field are Zariski
dense. Hence a rational section has \(D_Z(v)\ne0\) on every
component. These generic tests detect non-zero-divisors: the
components are reduced and integral, and the frame torus
separates distinct weights. \label{cor:covolume-component-transfer-14}

Expand this section in the short generators on each component.
When a coefficient on \(Z_0\) varies, the determinant on \(Z\)
becomes \(\det(A_Z+tB_Z)\), nonzero at the current value of \(t\).
Outside \(S Z_0\), \(B_Z=0\). Thus
\[
 \#\{t:\text{some determinant vanishes}\}
   \le \sum_{Z\in S Z_0}\deg_t\det(A_Z+tB_Z)\le |S|^2 .
\]
Choose a value in \(\{0,\ldots,|S|^2\}\) avoiding these roots,
and repeat for the other coefficients.

At each frame only one component contributes. With \(N_c\)
generators per component, Petersson bound \(\Lambda\), and comparison
factor \(\beta\), the triangle inequality gives
\[
                   \|v\|_{\sup}\le N_c |S|^2\beta\Lambda .
\] Here \(z_0=frL_{\rm sc}\),
\(8fr+10\le c_0\le22 f^2r^4\), and the chosen short
families and weight estimates give
\(\log N_c\le C_3 z_0^3\), \(\log\Lambda\le2C_3z_0^9\)
and \(\log\beta\le C_1' z_0^3\).
Thus the replacement by integral coefficients preserves the
order \(O(z_0^9)\) in the logarithm uniformly. Pullback by a
level automorphism is an isometry, and pointwise multiplication
followed by integration bounds every translate multiplier by
this same sup norm. \label{cor:covolume-component-transfer-15}
\end{proof}

\begin{construction}
\[
\begin{array}{lc}
\text{step on a nonterminal cover}&\text{upper bound for slope lost}\\ \hline
\text{character trace in the squared norm}&2b\\
\text{relative rank allowance }b/K & AU+5b\\
\text{descent of the coefficient lattice}&2b\\ \hline
\text{total beyond the deep bound }-M_i-AU&AU+9b
\end{array}
\]
\caption{The losses in passing from the shifted deep lattice to the base type. Here \(A=C_F+C_U\), and \(0<b\le\min(K,U)\) on the selected family.}\label{fig:component-trace-budgets}
\end{construction}

\paragraph{The definite endpoint.}
At the last lift \(q_m=0\), so the underlying definite quotient has finitely many components. Descent by traces between two positive-dimensional covers does not compute its degree. We compute the degree on each definite point and then pass to its type orbit. Put \(r=p_m\), let \(\Sigma\) be the chosen CM type, and set \(\ell=\log\Omega\). At \(\tau\in\Sigma\) write \(\alpha_\tau\) for the partition of the polynomial representation at the terminal weight, \(d_\tau=\dim S_{\alpha_\tau}(\mathbb C^r)\), and \(n\) for the even theta exponent. The total weights in this normalization are
\[
 \lambda_{\tau a}=\alpha_{\tau a}+n/2\quad(1\le a\le r),\qquad
 M_m=r\ell\sum_{\tau\in\Sigma}w_{m\tau}.
\]

\begin{auxlemma}[Evaluation and determinant of the definite Schur lattice]\label{lem:definite-schur-gram}
Let \(\alpha\) be a partition with at most \(r\) parts. Realize its dual Weyl lattice in flag polynomials on \(r\)-by-\(r\) matrices with the semistandard bideterminant basis of Proposition~\ref{prop:integral-terminal-branching}. Write \(d\) for its rank. At an orthonormal frame use the invariant norm with the extremal coefficient of norm one. If \(G_\alpha\) is the Gram matrix in the integral basis, then
\[
 0<\det G_\alpha\le d^d,\qquad
       |F(u)|\le\|F\|\quad (u\in U(r)).
\]
At the selected definite endpoint these determinants cost at most \(3U\) in the squared slope, summed over the CM type.
\end{auxlemma}
\begin{proof}
With Haar probability \(du\), Schur orthogonality identifies the squared norm of a flag polynomial with \(d\int_{U(r)}|F(u)|^2du\). Indeed these functions are matrix coefficients with a fixed unit extremal vector, as in Section~\ref{subsec:selected-column-metric}. Each integral basis polynomial is a product of minors. A minor of \(u\) is a matrix entry of \(\bigwedge^a u\) between coordinate unit vectors and has modulus at most one. Thus each diagonal Gram entry is at most \(d\). Positive definiteness and Hadamard's inequality give
\[
     0<\det G_\alpha\le\prod_{j=1}^d(G_\alpha)_{jj}\le d^d .
\]
For an orthonormal basis \(F_j\), the function \(\sum_j|F_j(u)|^2\) is constant by transitivity and unitary invariance. Its integral is \(d\cdot d^{-1}=1\). Cauchy--Schwarz gives the asserted evaluation bound. Tensoring at the CM labels preserves this operator bound, and tensoring by a Hermitian line introduces exactly its own norm.

To estimate the accumulated determinants take \(L=2L_{\rm sc}\). The terminal partition differs from the compact highest weight only by the scalar shift of all entries, or by taking the dual according to the frame convention. Both operations preserve the dimension. The selected weight bound in Lemma~\ref{lem:covolume-metric-factor} gives \(\alpha_i-\alpha_j\le2e^{2L}\) for \(i<j\). For \(L\ge2\), Weyl's formula gives
\[
 \log d_\tau
 =\sum_{i<j}\log\!\left(1+\frac{\alpha_{\tau i}-\alpha_{\tau j}}{j-i}\right)
 \le \tfrac32r^2L,\qquad
 \sum_{\tau\in\Sigma}\log d_\tau\le3fr^2 L_{\rm sc}\le3U .
\]
Here \(r^2\le b\) and \(m\log\log x\ge1\) on the selected run. This dimension estimate uses only the weight bound from that lemma.
\end{proof}

\begin{auxlemma}[Total displacement of the terminal weights]\label{lem:terminal-centred-weights}
On the selected size range, with the weights and period of the definite endpoint,
\[
 \left|\ell\sum_{\tau,a}(\lambda_{\tau a}-w_{m\tau})\right|
     \le |\ell|\sum_{\tau,a}|\lambda_{\tau a}-w_{m\tau}|\le 100U .
\]
\end{auxlemma}
\begin{proof}
Fix \(\tau\), put \(e=w_{m\tau}\) and use \(r=k+m-1\). The terminal weight formula for the first \(k-1\) coordinates has the same head parameters and rank shifts as in the Fock-degree calculation of Proposition~\ref{prop:selected-exponent-centering}. Number these coordinates starting at one, and denote their head parameters by \(0\le v_a\le S_c\). With \(t\) the block-length parameter and \(n=n_V\) even, they are
\[
 \lambda_{\tau a}
   =\left\lfloor v_a+e-(r-a)-\tfrac{t+1}{2}\right\rfloor
            +\tfrac{n_V}{2}\qquad(1\le a\le k-1).
\]
Thus each of these differences from \(e\) has absolute value at most \(S_c+r+t+n_V+2\), including the rounding error.

The remaining \(m\) coordinates use the node offsets at the same place. Before adding the rank shift their absolute values are at most the width \(2\eta_c\) of the offset window. The rank shift is at most \(2r\), as in the same Fock-degree calculation. Since \(0\le e\le3b\), \(\eta_c\le b\) and \(r\le b\), these coordinates obey
\[
   |\lambda_{\tau a}-e|\le 2\eta_c+2r+3b\le7b
                   \qquad(k\le a\le r).
\]

The selected size inequalities give \(rS_c\le16b\), \(r^2\le rn_V\le b\), \(t\le n_V\) and \(r\le b\). For the first inequality one may use \(r\le2k\) and the selected bound \(kS_c\le 8b\). The second follows from \(frn_VL_{\rm sc}\le b\). Consequently
\[
 \sum_a|\lambda_{\tau a}-e|
     \le r(S_c+r+t+n_V+2)+7mb\le21b+7mb .
\]
There are \(f\) labels and \(|\ell|\le3L_{\rm sc}\) in the selected CM data. Since \(m\ge1\), \(\log\log x\ge1\) eventually and \(U=fmbL_{\rm sc}\log\log x\), the right side summed with the period weight is at most \(84U\).
\end{proof}

\begin{definition}[Valuative frames on a definite component]\label{hyp:terminal-cm-fibre-comparison}
Use the geometric row and the level markings of
Construction~\ref{hyp:family-prescribed-realization}. At the terminal signature
\((r,0)\), its components over the splitting coefficient field \(K\) are
points. A test of integrality at a place \(v\) consists of a framed generic point
over a complete discrete valuation ring \(R\) above \(v\).
It is read on a semiabelian extension of its underlying polarized point.
Write \(M\) for its Lie lattice and \(\Theta\) for its rigidified principal
theta line. A finite frame \((e_{\tau,j},s)\) satisfies
\[
  \bigoplus_{j=1}^r R e_{\tau,j}
      =M\cap(M_{\operatorname{Frac}R})_\tau
         \quad(\tau\in\Sigma),\qquad Rs=H^0(A_R,\Theta).
\] The geometric lattice consists of the rational cusp forms
whose values belong to \(R\) in every such test. At infinity the test
uses the tangent metric \(\pi H\) of the polarization and the squared
cubical norm of the generator integrated against Haar probability.
These conventions fix both a lattice and a metric.

Write \(C_\tau,T\) for the coefficient determinant and theta lines,
\(c_{\tau,Z},t_Z\) for their unsquared degrees and set
\[
 \Delta_Z=\max\!\left(0,\sum_\tau c_{\tau,Z}-2t_Z\right),
 \qquad C_D=105 .
\]
\end{definition}

At \(q_m=0\) the conjugate determinant has rank zero. The coefficient
and flat determinant generators transform inversely. Their integral
eigenframes and polarized metrics in
Definition~\ref{hyp:flat-determinant-frames} are dual. Applying
Proposition~\ref{prop:deep-flat-line-degree} gives
\[
       c_{\tau,Z}=-\widehat{\deg}_{\rm flat}(\mathcal L_\tau)
                  =-r\ell/2 .
\]

\begin{auxlemma}[The theta comparison on a principal CM fibre]\label{lem:terminal-principal-theta-comparison}
After finite base extension to a field \(K\) containing all conjugates of \(E\), take a semistable Moret--Bailly model of the principally polarized CM point, with rigidified symmetric cubical theta line \(\Theta\), integral \(\mathcal O_E\)-action and finite flat polarization kernel for \(\Theta^2\). Write \(M\) for its integral Lie lattice with the polarization metric \(\pi H_\sigma\), and \(B\) for the direct image of \(\Theta\) with the cubical probability \(L^2\) metric. Put \(M_\tau=M\cap(M_K)_\tau\), \(C_\tau^0=(\det M_\tau)^\vee\), and \(T^0=B^\vee\).

Let \(C_\tau,T\) be the actual metrized coefficient lines of degrees \(c_{\tau,Z},t_Z\). Define the line measuring the difference of frames and its degree by
\[
 \mathcal E_Z=\left(\bigotimes_{\tau\in\Sigma}\bigl(C_\tau\otimes(C_\tau^0)^\vee\bigr)\right)
             \otimes(T^0\otimes T^\vee)^{\otimes2},
     \qquad d_Z=\widehat{\deg}\mathcal E_Z .
\]
For the definite signature, let \(g=fr\) and let \(I_Z\) be the normalized logarithmic order of \(M/(\bigoplus_{\tau\in\Sigma}M_\tau)\). Then
\begin{equation}\label{eq:terminal-theta-frame-defect}
 \sum_{\tau\in\Sigma}c_{\tau,Z}-2t_Z
      = I_Z-\tfrac g2\log(2\pi)+d_Z,\qquad
      0\le I_Z\le g\log\operatorname{rd}(E).
\end{equation}
In particular, on matched integral and metric frames \(d_Z=0\), and the whole theta contribution to \eqref{eq:terminal-cm-loss} satisfies \(n\Delta_Z\le2U\).
\end{auxlemma}
\begin{proof}
For a principal theta line \(h^0(\Theta)=1\). In the notation \(h_*=-\widehat{\deg}\det M\) of Section~\ref{subsec:flag-point-heights}, the direct-image formula at the first power of \(\Theta\) is
\[
 \widehat{\deg} B=-\tfrac12 h_*-\tfrac g4\log(2\pi),\qquad
 2\widehat{\deg} T^0=-\widehat{\deg}\det M+\tfrac g2\log(2\pi).
\]
This is the rank-one case of Bost's formula for the cubical direct image, also expressed in the frame convention in \eqref{eq:slice-cm-theta-degree}. \cite[Gaudron, 2019, \S2.3.2]{bost-mb-flag}
The exponent \(n\) of the automorphic weight is applied to this one coefficient line. A section evaluated on a theta generator of squared norm \(N_\sigma\) has coefficient-frame factor \(N_\sigma^{-1/2}\). Thus the standard coefficient line is the dual of the direct image; taking \(H^0(\Theta^n)\) would change its rank.

Since \(q_m=0\), the conjugate half contributes no generic Lie eigenspaces. The polarization pairs the \(\mathcal O_E\)-action with its complex conjugate under the adjoint, so the distinct complex eigenspaces are orthogonal in the Hermitian metric. The remaining determinant discrepancy is the finite index:
\[
 \widehat{\deg}\det M=\sum_{\tau\in\Sigma}\widehat{\deg}\det M_\tau+I_Z .
\]

At finite places the integral decomposition can fail even at this signature. Let \((u_j)_j\) be an integral basis of \(\mathcal O_E\) and \((u_j^*)_j\) its trace-dual basis. The projector \(e_\tau=\sum_j u_j\otimes\tau(u_j^*)\) belongs to \(E\otimes K\); evaluation at embedding \(\rho\) gives \(\delta_{\rho\tau}\) by trace duality. Since \(u_j^*\in\mathfrak d_E^{-1}\), multiplication by \(\tau(\mathfrak d_E)\) clears the denominators of \(e_\tau\). At a finite place \(v\) of \(K\) it follows that
\[
 \bigoplus_\tau M_{\tau,v}\subset M_v
      \subset\bigoplus_\tau \tau(\mathfrak d_E)^{-1}M_{\tau,v},\qquad
 \operatorname{length}_{\mathcal O_{K,v}}(M_v/\textstyle\bigoplus_\tau M_{\tau,v})
     \le r\sum_{\tau\in\Sigma}v(\tau\mathfrak d_E).
\]
The normalized sum of \(v(\tau\mathfrak d_E)\log\mathrm Nv\) is \(\log\operatorname{rd}(E)\) for each \(\tau\). This proves the index bound. Subtracting twice the degree of \(T^0\) from the sum of the degrees of \(C_\tau^0\) and adding \(d_Z\) gives \eqref{eq:terminal-theta-frame-defect}.

The comparison can be tested on the actual fibres. For \(0\ne e\in\mathcal E_Z\otimes K\) write its finite lattice as \(J_v e\) locally and its length as \(q_\sigma(e)\) at each complex embedding. Then
\[
 d_Z=-\frac{\sum_v v(J_v)\log\mathrm Nv+\sum_\sigma\log q_\sigma(e)}
                    {[K:\mathbb Q]},\qquad
 n\Delta_Z\le n g\log\operatorname{rd}(E)+n\max(0,d_Z).
\]
The right side is independent of the rational frame by the product formula. On the selected family \(n=n_V\), \(frn_V L_{\rm sc}\le b\), \(\log\operatorname{rd}(E)\le2L_{\rm sc}\), and \(b\le U\). This proves the stated allowance.
\end{proof}

\begin{auxlemma}[The integral lattice at the selected definite points]
\label{lem:terminal-valuative-realization}
Use the geometric lattice of Definition~\ref{hyp:terminal-cm-fibre-comparison}.
At each terminal point the integral polynomial fibre in the tangent and
theta frames contributes its full rank to the geometric cusp space.
The selected level subgroup acts freely on the points and preserves
the corresponding lattices and norms. In the notation of
Lemma~\ref{lem:terminal-principal-theta-comparison}, one has
\begin{equation}\label{eq:terminal-cm-loss}
 C_\tau=C_\tau^0,\qquad T=T^0,\qquad d_Z=0,\qquad n\Delta_Z\le2U .
\end{equation}
The equalities of lines are isometries with their integral structures.
\end{auxlemma}
\begin{proof}
\proofstep{Retain one marking prime to the residue characteristic.}
The terminal abelian variety \(A\) has dimension \(fr\), is principally
polarized, and has an \(\mathcal O_E\)-action. Its Lie algebra has type
\(r\Phi\) for the selected CM type \(\Phi\). In characteristic zero
the idempotent decomposition shows that the CM-type ideal
\[
 J_\Phi=\ker\!\left(\mathcal O_E\otimes\mathcal O_{E^*}
    \longrightarrow \prod_{\tau\in\Phi}\overline{\mathbb Q},
              \ a\otimes b\longmapsto(\tau(a)b)_\tau\right)
\]
annihilates \(\operatorname{Lie}A\); here \(E^*\) is the reflex field
in the fixed algebraic closure. The integral moduli stack of these
principally polarized triples with this ideal condition is proper and
smooth of relative dimension zero over \(\mathcal O_{E^*}\).
\cite[Theorem~35]{amir-khosravi-serre}

Fix a place \(v\) of \(K\). Forget the part of the level above its residue
characteristic and retain the principal marking at one of the two
distinct split primes \(g_0,g_1\) prime to that characteristic.
Adding this marking gives a finite \'etale cover of the stack over the
local base: the relevant half of the \(g_i\)-torsion is finite \'etale,
and its bases form a finite \'etale sheaf. For a marked split factor its geometric fibre is a torsor under
\[
 \operatorname{Aut}_{\mathbb F_{g_i}}
     \bigl(A[\mathfrak g](\overline{k})\bigr)
           \simeq\mathrm{GL}_r(\mathbb F_{g_i}),\qquad \mathfrak g\mid g_i .
\]

This cover has trivial inertia. Indeed a polarized automorphism \(u\) commuting with \(E\) has finite
order. On the marked \(\mathfrak g\)-part of the Tate lattice,
\(\mathfrak g\mid g_i\), it acts trivially modulo \(g_i\).
The principal congruence kernel in \(\mathrm{GL}_r(\mathbb Z_{g_i})\)
is torsion-free since \(g_i\) is odd. The action on this part is faithful:
\[
 \operatorname{End}^0_{E}(A)\ \lhook\joinrel\longrightarrow\
 \operatorname{End}^0_{E}(A)\otimes_E E_{\mathfrak g}
 \ \lhook\joinrel\longrightarrow\
 \operatorname{End}_{E_{\mathfrak g}}(V_{\mathfrak g}A).
\]
For the second injection take the \(\mathfrak g\)-factor of the
faithful rational Tate-module realization of endomorphisms at
\(g_i\), which is invertible in the residue field.
This applies on special as well as generic fibres.

The marked cover is therefore a proper algebraic space, and its given
\(K_v\)-point extends over \(\mathcal O_{K_v}\). Write \(\mathcal M_{\Phi,\mathfrak g}\) for this space and \(\mathcal A\)
for its universal abelian scheme. For a testing trait \(R\) above \(v\)
the good model is the pullback
\[
  \operatorname{Spec} R\longrightarrow\operatorname{Spec}\mathcal O_{K_v}
      \xrightarrow{\,\widetilde A\,}\mathcal M_{\Phi,\mathfrak g},
    \qquad A_R=(\widetilde A^*\mathcal A)_R .
\]
All evaluations retain the original generic level.

\proofstep{Identify all integral and metric frames.}
Any semiabelian extension in a test is canonically the base change
of this good model. Indeed the identity of the generic fibres and its
inverse both extend over a normal trait by the extension theorem for
homomorphisms of semiabelian schemes. Uniqueness makes the extensions inverse.
\cite[Chapter~I, Proposition~2.7]{faltings-chai}

On this abelian scheme the rigidified theta line extends uniquely
up to the isomorphism inducing the chosen rigidification. One may
extend a Cartier divisor from the generic fibre by closure on the regular
scheme, and rigidify along zero. For two extensions \(\Theta_1,\Theta_2\) the ratio of a fixed
generic isomorphism is given by a vertical divisor. The special fibre
\(A_k\) is integral. If \(e\) denotes the zero section, the divisor and
its pullback therefore have the form
\[
      D=m[A_k],\qquad e^*D=m[\operatorname{Spec} k].
\]

Compatibility with the given rigidifications forces \(m=0\). The cubical structure is the canonical
one on the rigidified extension. The polarization remains principal,
and formation of the rank-one module of sections commutes with base
change (the higher cohomology of an ample line on an abelian variety
vanishes). Writing \(M=\operatorname{Lie} A_R\) for the good lattice,
the two tests are consequently
\[
    M_\tau=\bigcap_{a\in\mathcal O_E}\ker(a-\tau(a)\operatorname{id}_M),
       \qquad B=H^0(A_R,\Theta).
\]
The first kernel is the intersection of the generic eigenspace with
\(M\), and commutes with flat extension of traits. The coefficient
coordinates take the two duals \((\det M_\tau)^\vee\) and \(B^\vee\).
The metrics in Definition~\ref{hyp:terminal-cm-fibre-comparison} are exactly
the dual metrics used in Lemma~\ref{lem:terminal-principal-theta-comparison}.
This proves \eqref{eq:terminal-cm-loss}, including at ramified primes.

\proofstep{Take sections supported at a point.}
In the split definite quotient a rule of the given weight over a point
is an element of the corresponding polynomial representation tensored
with its theta line power. This follows by evaluating in a frame and
then using the weight transformation law in all other frames.
Let \(X_K\) denote the finite split quotient and \(V_Z\) its coefficient
space at \(Z\). Each point is open and closed and the definite group
has no cuspidal boundary. The extension map is
\[
 \iota_Z:V_Z\longrightarrow\bigoplus_{Z'\in X_K}V_{Z'},\qquad
    (\iota_Z v)_{Z'}=
       \begin{cases}v&Z'=Z,\\0&Z'\ne Z.\end{cases}
\]

On its own point an integral polynomial
has integral values in all the frames just identified, by the integral
change-of-basis action on the Schur lattice. On the other points all
tests have value zero. Thus extension by zero supplies an integral map
whose image spans the entire rational fibre. This uses precisely the
valuative lattice, even when two marked generic points have the same
specialization on a model retaining the full level.

Write \(P\) for the selected type subgroup of the two-prime level
quotient. If \(s\in P\) fixes \(Z\), its action on the marking is realized
by a polarized \(\mathcal O_E\)-linear automorphism \(u\) of \(A_Z\):
\[
    sZ=Z\quad\Longrightarrow\quad
       u\circ\eta_i=\eta_i\circ s_i \qquad(i=0,1),
\]
where \(\eta_i\) denotes the retained part of the principal marking
and \(s_i\) its matrix at \(g_i\). 

The matrices are unipotent on a split half at both primes, as in the orbit
calculation in this chapter. On the part at \(g_i\) the automorphism lies in the inverse image of the cyclic unipotent subgroup generated
by its reduction. Write \(u_i\) for its action on this Tate lattice. The reduction kernel
is pro-\(g_i\) and the cyclic image is a \(g_i\)-group. Hence, by the
faithfulness already proved,
\[
    \operatorname{ord}(u)=\operatorname{ord}(u_i)=g_i^{a_i}
       \quad(i=0,1),\qquad (g_0,g_1)=1 .
\]
Thus \(u=1\). Faithfulness of the marking gives the identity in the type group. Transport
changes only the marking and carries both the valuative frames and the
polarized Hermitian norm to those of the translated point.
\end{proof}

\begin{auxlemma}[The terminal fibre gives the type slope]\label{lem:terminal-cm-type-slope}
For the geometric lattice of Lemma~\ref{lem:terminal-valuative-realization}, the type slope at the terminal lift satisfies
\[
 \mu_m\ge-M_m-C_D U
       =-r\log\Omega\sum_{\tau\in\Sigma}w_{m\tau}-C_D U .
\]
\end{auxlemma}
\begin{proof}
For a partition \(\alpha\) put \(s=|\alpha|\), \(d=\dim S_\alpha\). The determinant of the polynomial representation is a character of \(\mathrm{GL}_r\). Restricting to diagonal matrices, permutation symmetry gives equal exponents on the \(r\) entries, while a scalar \(zI\) acts as \(z^s I_d\). Therefore
\[
 \det S_\alpha(g)=(\det g)^{sd/r}.
\]
The exponents can also be read directly from the diagonal weights. The determinant of a tensor product satisfies \(\det(A\otimes B)=(\det A)^{\dim B}(\det B)^{\dim A}\), so the normalized determinant degrees add. The integral and metric comparisons give for the squared fibre slope
\[
 \mu_Z\ge 2\sum_\tau\frac{|\alpha_\tau|}{r}c_{\tau,Z}
                       -\sum_\tau\log d_\tau+2nt_Z.
\]
Since \(2t_Z\ge\sum_\tau c_{\tau,Z}-\Delta_Z\), substitution gives
\[
 \mu_Z\ge-\ell\sum_{\tau,a}\lambda_{\tau a}
                      -\sum_\tau\log d_\tau-n\Delta_Z
             \ge-M_m-C_D U
\]
by \eqref{eq:terminal-cm-loss} and Lemmas~\ref{lem:definite-schur-gram} and~\ref{lem:terminal-centred-weights}. The sign in the weight term is fixed by the common scalar \(-\ell\); no absolute value is needed before summing.

For a representative vector \(v\) on \(Z\), the section
\[
             v\longmapsto\sum_{g}\chi(g)^{-1}g v
\]
has type \(\chi\). Its restriction to \(Z\) recovers \(v\), by freeness of the orbit. Conversely, a rational section of this type is determined by its values at one representative per orbit. Taking the direct sum therefore gives an integral map of full rank into the type lattice. The finite index of its image can only increase the degree of the target. For each orbit let \(m_Z\) be its mass in quotient probability measure. Only one translate can contribute at a point of that orbit. The tensor evaluation bound of Lemma~\ref{lem:definite-schur-gram}, invariance of the norm and \(|\chi(g)|=1\) give
\[
 \left\|\sum_{Z,g}\chi(g)^{-1}g v_Z\right\|_\sigma^2
    \le\sum_Z m_Z\|v_Z\|_{Z,\sigma}^2
    \le\sum_Z\|v_Z\|_{Z,\sigma}^2 ,
 \qquad m_Z\ge0,\quad\sum_Zm_Z=1 .
\]
The induced map on highest exterior powers is thus contractive at each embedding. Hence \(\mu_m\) is at least the rank-weighted mean of the \(\mu_Z\).
\end{proof}

\begin{proposition}[Descent by the character traces]\label{prop:covolume-trace-descent}
Write \(\mu_i\) for the slope of the base type on the nonterminal cover. Then
\[
 \mu_i\ge-M_i-\{2(C_F+C_U)+9\}U.
\]
 At the terminal lift Lemma~\ref{lem:terminal-cm-type-slope} gives \(\mu_m\ge-M_m-C_D U\).
\end{proposition}
\begin{proof}
For \(h\in S\) form the character trace
\[
 \phi_h(f)=\sum_{g\in S}\chi(g)^{-1} f^g v^{hg}.
\]
The level action preserves integrality and cuspidality, and the
even weights give the required target weight. Reindexing shows
that each trace transforms by \(\chi\). Put \(A=(v^{hg})_{h,g}\). Vanishing of the traces gives the system
\[
 A\bigl(\chi(g)^{-1} f^g\bigr)_{g\in S}=0,\qquad
 \operatorname{adj}(A) A=(\det A) I .
\]
The determinant is a non-zero-divisor, hence \(f=0\). In the
opposite direction, translate a relation \(\sum_h f_hv^h=0\)
on type vectors by each \(g\). Cancelling \(\chi(g)\) gives the
transposed system, hence injectivity and
\(|S|\operatorname{rk} L_{\rm deep}\le\operatorname{rk} L_+\).
\label{cor:covolume-component-transfer-16}

At each embedding, with \(B_v\) a bound for multiplication by
each translate of \(v\), isometry of the action and \(|\chi(g)|=1\)
give
\[
 \|\phi_h(f)\|\le |S|B_v\|f\|,\qquad
 \sum_{h\in S}\|\phi_h(f)\|^2\le |S|^3 B_v^2\|f\|^2 .
\]
The uniform height and level estimates for the chosen section
give \(\log B_v\le C z_0^9\), \(\log |S|\le45z_0^2\),
where \(z_0=fr L_{\rm sc}\). Eventually \(z_0\ge2C+200\)
and \(z_0^{12}\le b\). Thus its squared-norm factor satisfies
\[
 \log(|S|^3 B_v^2)\le135z_0^2+2Cz_0^9\le2z_0^{12}\le2b.
\] Write \(K=P\kappa_i\) for the size bound on
\(|M_i|\) and the log size of a spanning family in the deep
type lattice. The selected
family has \(0<b\le\min(K,U)\). On the deep cover the
negative-shift lattice has squared-norm slope at least
\(-M_i-(C_F+C_U)U\).

The shifted-rank comparison on the same cover gives
\[
 \operatorname{rk}L_+\le(1+b/K)\operatorname{rk}L_-,
 \qquad
 \frac{\operatorname{rk}L_+-\operatorname{rk}L_-}
      {\operatorname{rk}L_-}\le b/K .
\]
The same relative rank allowance applies after summing over the
trace targets. Writing
\(v=-M_i-(C_F+C_U)U-2b\), the resulting slope is at least
\(v-(b/K)(|v|+2K)\). Descent through the coefficient ring
costs another \(2b\). To combine the rank and descent costs put \(A=C_F+C_U\). Since
\(|M_i|\le K\), \(b/K\le1\), and \(b\le U\),
\[
 \frac bK (|v|+2K)
   \le b+AU+2b+2b=AU+5b .
\]
The rank loss and the two metric losses of \(2b\) each thus give
a base type slope of at least
\[
                    -M_i-\{2(C_F+C_U)+9\}U .
\]
This uses the trace maps, minimum-slope, degree, share and
Hilbert--Samuel bounds on the stated integral models.
At the terminal lift Lemma~\ref{lem:terminal-cm-type-slope} applies instead. \label{cor:covolume-component-transfer-17}
\end{proof}
\paragraph{Norm used for the arithmetic row.}
The bound so far uses the geometric metric. The arithmetic row carries an adelic norm,
so we must compare the two on the type lattice just estimated. Fix its scalar-model embedding, the local
compact lowest types of dimensions \(d_{\tau,i}\), and a
fundamental domain \(\mathcal F_i\) for the corresponding
adelic group modulo its rational points. Use the Haar measure
\(\mu_i\) with the real density and finite-level normalization
of the scalar model. For a scalar section \(f\), write
\(\|f\|_{\rm P}=a_i\|f\|_{{\rm Pet},1}\), where
\(\|f\|_{{\rm Pet},1}\) is the Petersson norm with quotient
probability measure and \(a_i>0\) is the norm factor of the
specified Lie frame in this weight. The adelic scalar evaluation contains this same factor.
\begin{lemma}[Normalization of the type metric]\label{lem:covolume-type-metric}
The squared multiplicity norm is \(R_i\|f\|_{\rm P}^2\), where
\[
            R_i=\mu_i(\mathcal F_i)\prod_\tau d_{\tau,i}.
\]
\end{lemma}
\begin{proof}
Evaluate on the tensor of the prescribed unit
lowest-type vectors. The scalar attached to the form is this value or its complex conjugate. Multiplying it by \(a_i^{-1}\)
gives a function on the probability quotient whose \(L^2\)
norm is \(\|f\|_{{\rm Pet},1}\), by the scalar Petersson
comparison on the chosen model. Schur orthogonality gives
the squared evaluation integral as the squared multiplicity
norm divided by \(\prod_\tau d_{\tau,i}\). Let \(d\mu_{i,1}\) be quotient probability measure. The measure
conversion in this calculation is
\[
        d\mu_i=\mu_i(\mathcal F_i)\,d\mu_{i,1}.
\]
It multiplies the squared scalar integral by precisely this volume. Cuspidality gives
the integrability used here. For the conjugate coefficient
embedding the canonical good embedding is the same and
complex conjugation preserves the squared norm of this scalar. This gives the multiplier \(\sqrt{R_i}\) at every
embedding of the actual row field.
\label{cor:covolume-component-transfer-18}

For a finer compact open \(K\subset K_0\), keeping the real
density and giving the chosen finite level volume one
multiplies the Haar measure by \([K_0:K]\). The quotient
probability measures agree under passage to the coarser
level: on the full adelic quotient a constant change of
Haar cancels against total mass, and the pushforwards are
the probability measures on the respective level quotients.
Writing \(p:Y_K\to Y_{K_0}\) for the projection and
\(\mu_{K,1},\mu_{K_0,1}\) for the quotient probability measures gives
\[
 p_*\mu_{K,1}=\mu_{K_0,1},\qquad
 \int_{Y_K}|p^*s|^2\,d\mu_{K,1}
    =\int_{Y_{K_0}}|s|^2\,d\mu_{K_0,1}.
\] The multiplicity comparison at
the finer level uses that level's own Haar volume. Apply this
calculation to each run's marked subgroup and independently
transport its Lie frame and rational coefficient reading.
The later finite coordinate factor \(n_P^{-1}\) in squared
norm comes from the row frame isometry.

\end{proof}

\begin{lemma}[Size of the metric factor]\label{lem:covolume-metric-factor}
Use the marked-level factor \(R_i\) of Lemma~\ref{lem:covolume-type-metric}.
Let \(L\ge2\) with \(r+3\le frn_V\le e^L\), normalized absolute
weight sum at most \(e^{2L}\), and
\[
 \log(g_0g_1)\le30L,\qquad
 \log\operatorname{rd}_E+|\RePart(L'/L)(1,\chi_K)|\le2L .
\]
Then at each marked level
\[
 -fr^2\log64\le\log R_i\le122fr^2 L .
\]
In particular \(|\log R_i|\le244U\) for \(L=2L_{\rm sc}\) eventually.
\end{lemma}
\begin{proof}
Let
\(K\subset K_0\) be the selected level and full level and
\(I=[K_0:K]\). Cover a fundamental set by translates
indexed by the finite lattice classes at full level and
then by the \(I\) level cosets. In real unitary coordinates
the determinant torus and the special unitary factor give
an upper density and box factor \((1000fr^3)^{fr^2}\).
The finite translates are represented by self-dual lattices.

Grouping them by their determinant class, with the
congruence-three marking, bounds the sum of the special
unitary covolumes by
\[
 T_E=h_E\,3^{2f+2fr^2}\,
       \operatorname{rd}_E^{\,f(r^2-1)/2}
       \left(\sqrt r\prod_{a=2}^r\zeta(a)\right)^f .
\]
Here \(h_E\) is the class number. This is the special
unitary Tamagawa comparison for the unramified CM extension
and the specified self-dual integral lattice with the tangent
Riemannian Haar measure. The determinant class count and the
index of the level-three marking account for the class number
and powers of three. Subadditivity on the fundamental set
therefore gives \(\mu_i(\mathcal F_i)\le I(1000fr^3)^{fr^2}T_E\).
\label{cor:covolume-component-transfer-19}

For the reverse inequality take special unitary double-coset
representatives at full level. A determinant sector of the real
domain times special unitary fundamental sets packs into the
unitary quotient with lower factor \(32^{-fr^2}\). Overlaps
can occur only by a rational unitary transformation with
integral determinant. In the chosen sector it has determinant
one and the special unitary representatives then coincide.
The sum of these special unitary covolumes is
\[
 |\operatorname{disc}(E^+)|^{(r^2-1)/2}
             (\sqrt r)^f/P_{\rm fin} .
\]

In the Tamagawa formula the tangent normalization contributes
the index of the chosen integral Lie basis times
\((\sqrt r)^f\), and the inverse index from the rational
volume form cancels it. For the unramified local factors,
\[
 0<P_{\rm fin}\le\left(\prod_{a=2}^r\zeta(a)\right)^f\le2^{fr}.
\]
Inserting this upper bound into the denominator bounds the sum
of covolumes below by \(2^{-fr^2}\).
Passage to the smaller level increases mass in the
level-normalized Haar measure. Since \(d_{\tau,i}\ge1\),
this gives \(\log R_i\ge-fr^2\log64\).
\label{cor:covolume-component-transfer-20}

For the upper comparison take the size parameter \(L\ge2\)
with \(frn_V\le e^L\), normalized absolute weight sum at
most \(e^{2L}\), \(\log(g_0g_1)\le30L\) and
\(\log\operatorname{rd}_E+|\RePart(L'/L)(1,\chi_K)|\le2L\).
Bounding each finite reduction group by all matrices over its
residue field gives the index budget
\[
 I\le(g_0g_1)^{2fr^2},\qquad
 \log I\le2fr^2\log(g_0g_1)\le60fr^2L .
\]
The discriminant and class-number bound in the preceding
upper estimate gives
\(\log\mu_i(\mathcal F_i)\le fr^2\log(r+3)+119fr^2L\).

The block Weyl dimension product bounds
\(\log d_{\tau,i}\le 3r^2 L/2\) on each of the \(f\) labels.
The selected ranks give \(r+3\le frn_V\), so
\(\log(r+3)\le L\). In the formula for the squared metric multiplier this yields
\[
 \log R_i=\log\mu_i(\mathcal F_i)+\sum_\tau\log d_{\tau,i}
    \le (1+119+3/2)fr^2L\le122fr^2L .
\] With \(L=2L_{\rm sc}\)
on the chosen run, \(r^2\le b\) and
\(U=fmbL_{\rm sc}\log\log x\ge fr^2 L_{\rm sc}\).

The two estimates imply the asserted absolute bound for \(R_i\).
\end{proof}

\begin{lemma}[Integral change of field for the type]\label{lem:covolume-integral-descent}
The integral type lattice is the scalar extension of the
smaller-field lattice. With \(\mu_{\rm row}\) its squared-norm slope under the row metric and \(\mu_{\rm small}\) the smaller-field slope,
\[
 \mu_{\rm row}\ge\mu_{\rm small}-244U.
\]

\end{lemma}
\begin{proof}
In a rational basis of the cusp model the type subspace is
the simultaneous kernel of the row actions minus their prescribed
eigenvalues; the latter belong to the smaller field. A finite
subsystem cuts out the same kernel. If \(M\) is the integral
cusp lattice and \(T\) this rational type subspace, its integral
part is \(\ker(M\to M_{\mathbb Q}/T)\), where \(M_{\mathbb Q}\)
denotes its span over the smaller coefficient field. The larger integer ring is torsion-free
over the smaller Dedekind ring, hence flat. Tensoring this kernel
therefore gives the intersection defining the type over the
larger field. At each embedding the norm multiplier is
\(\sqrt{R_i}\le e^{122U}\); in the squared degree the change per
rank is at least \(-244U\). \label{cor:covolume-component-transfer-21}
\end{proof}

\paragraph{The maximal slope of the full cusp lattice.}
\label{par:parent-full-cusp-slope}
The full cusp lattice and its row have different roles in this field comparison. Let \(K_c\)
be the field of the geometric cusp model at the base level,
and \(K\) the field carrying the row. The size estimate for the
recursive class sum bounds \(\log\operatorname{rd}(K_c)\) by
\(217(frL)^2\), since the field of the core is the coefficient
field in that estimate. We first work over \(K_c\), which is
totally complex. Write \(\|\cdot\|_{\rm Pet}\) for the adelic probability
Petersson norm multiplied by the Lie-frame factor \(a_i\)
above. Use the same factor on sup norms. It cancels in their
comparison. The integral lattice
has full rank there. A nonzero good component is supplied by the
base realization of rank \(D_0\). Thus every application of a
maximal-slope or a logarithmic rank bound below has positive
ambient rank. \label{cor:parent-cusp-slope-1}

\begin{established}[Sizes on the completed parent]\label{hyp:parent-cusp-size}
Use the geometric field \(K_c\), row field \(K\), and normalized
Petersson metric in the common primitive frame. Choose a fixed exponent \(a\) with \(frn_V L\le X=x^a\le b\) on the late family, where
\(L=\log b\ge10\), \(f\ge8\) and \(r\ge2\).
Write \(s\) for the sum of the absolute frame powers, \(t=frn_V\)
and \(w=1+s+t\). Summing the dominant weight bounds on the
two signature blocks gives \(s\le8frb\), hence
\(w\le10frb\) and \(\log w\le4L\).
With \(\ell=2L\), the two-prime level \(M=g_0g_1\), the
quadratic discriminant parameter \(d\) and \(d_E=|\operatorname{disc}E|\)
satisfy
\[
 \log(fr)\le\ell,\quad \log M\le30\ell,\quad
 \log d_E\le4f\ell,\quad \log d\le4\ell,\quad \log w\le2\ell.
\]
The marked CM period satisfies \(|\log\Omega|\le3L\).
\end{established}
\begin{proof}
The ranks of Lemma~\ref{lem:window-parameter-choice} and the number of columns in the completed point grow at most polynomially in \(x\); \(b=\exp(x^\epsilon)\) eventually dominates each fixed power of \(x\). The upper box bound \(2b\) per entry, applied to the two frame lists and their determinant shifts, gives \(s\le8frb\). On the selected family \(frn_VL\le x^a\le b\), hence \(w\le10frb\) and \(\log w\le4L\) for large \(x\). For the same field and two primes the calculation of Lemma~\ref{lem:deep-cover-field-size} gives \(\log g_j\le9L+1\), and its bound for the root discriminant gives \(\log |\operatorname{disc} E|=2f\log\operatorname{rd}(E)\le 4f\ell\). The two-prime product and the marked discriminant then give the remaining logarithmic bounds. The period bound is the CM period estimate on the retained family used in Lemma~\ref{lem:terminal-centred-weights}.
\end{proof}

\begin{proposition}[A uniform upper bound for the parent maximal slope]
\label{prop:parent-cusp-maximal-slope}
For each lift of the completed parent with
Lemma~\ref{hyp:parent-cusp-size}, the ambient integral cusp
lattice satisfies
\[
  \mu_{\max}(\mathrm{lat}_{\rm all})\le x^A b
\]
for a fixed \(A\), uniformly in the row field. The maximal slope
uses the normalized Petersson metric and the actual ambient
lattice in the primitive frame.
\end{proposition}

\begin{proof}
\proofstep{The lower norm bound for an integral cusp form.}
For a nonzero integral form, apply the CM reading
\eqref{eq:slice-cm-point-norm} with these powers. The cusp
flow at level \(M\) uses the thickness
\(e_0=\min(1,\pi/(2M\sqrt d(fw+2)\log2))\).
The CM approximation there gives a frame bound
\(\Lambda=R^{2q+1}\), where \(q\le r\) is the negative signature
on this lift. To use the negative block, approximate its \(q\)
axes in the signed Hermitian space and take the orthogonal
complement. The graph bound is invariant under passage to its
adjoint. The Gram determinant which clears the complementary
projection then has size at most \((2A)^{2q}\) in the notation
of the CM reading. This gives
\[
 R=\tfrac{1600}{3}r\sqrt{fw}\,(fr)^3d_E/\sqrt{e_0}.
\]
Indeed \(-\log e_0\le\log(2M\sqrt d f w)\).

Taking logarithms in the formula for \(R\) gives
\(\log R\le(4f+32)\ell\). The nonzero integral CM
coefficient, with its dual letter metrics, now gives
\(\operatorname{mean}_\sigma\log\|F\|_{\sup,\sigma}\ge-\lambda_s\),
where
\[
 \lambda_s\le
 \bigl((1+|\log\Omega|)/2+\ell+
                  (2q+1)(6f+48)\ell\bigr)(s+t).
\]
Here the value line contributes half the logarithm of
\(\Lambda\) per letter, and reading the frame coefficient costs
at most \(s\log(r\Lambda)\). \label{cor:parent-cusp-slope-2}

\proofstep{Returning to the Petersson norm.}
For the passage back to the Petersson norm, let \(B\) be the
bound on squared point evaluation. In a flag chart of complex
dimension \(D=fr(r-1)/2\) take \(h=(200r^3(w+f))^{-1}\).
On the part with larger weights the
holomorphic submean inequality contributes
\(8D!/(\pi h^2)^D\), including the variation of the frame norm
and the choice of a point above half the supremum. Overlap
of the charts is bounded by the number of integral transition
matrices moving the frame on the subspace with larger weights into a fixed ball. Packing their
entries as in the evaluation-net argument bounds this
multiplicity by
\[
 \mathcal M(e_0)=(1+2(2fr)^2d_E/e_0)^{2fr^2}.
\]
Integration on the
probability space introduces the upper level mass \(V_M\)
from the measure comparison. Thus one can take
\[
 B=\frac{8D!}{(\pi h^2)^D}\mathcal M(e_0)V_M,\qquad
             \|F\|_{\sup,\sigma}^2\le B\|F\|_{\rm Pet,\sigma}^2 .
\]

For clarity, the chart factor has logarithm at most
\(3\log2+fr^2(12\ell+16\log2)/2\). In the packing bound use
\(T=3fwM\sqrt d\ge e_0^{-1}\). Its logarithm is at most
\(2\log2+35\ell\). The log multiplicity is at most
\(2fr^2(\log16+2\log(fr)+\log d_E+\log T)\).
The log mass is at most \(130fr^2\ell\).
Adding gives \[
 0\le\log B\le250 f^2r^2\ell.
\]
These bounds use the equivariant frame character at the level,
the even central weight, and the same Haar probability as the norm.
\label{cor:parent-cusp-slope-3}

\proofstep{The rank bound using the same finite net.}
The dimension bound uses this identical \(B\). At an embedding
of \(K_c\), a rationally independent frame of forms stays
independent in the space of holomorphic readings by the
descent in paragraph~\ref{par:slice-evaluation-net}. If \(F_1,\ldots,F_d\) is an orthonormal basis for the
Petersson probability measure \(\mu\), the evaluation functional
at \(z\) has squared norm \(\sum_i\|F_i(z)\|^2\). The evaluation
inequality applies to every vector in their span, so
\[
 d=\int\sum_{i=1}^d\|F_i(z)\|^2\,d\mu(z)
       \le\int B\,d\mu(z)=B.
\]
This is the same \(B\) as in the evaluation bound.

Now in any nonzero integral sublattice of rank \(u\) and
squared-norm slope \(\mu\), the Minkowski bound used for the
short section gives a nonzero form with
\[
 2\operatorname{mean}_{\sigma}\log\|F\|_{\rm Pet,\sigma}
      \le -\mu+\log u+\log\operatorname{rd}(K_c).
\]
Indeed the Euclidean ball of radius one in each complex rank
\(u\) contains the product of coordinate discs of radii
\(1/\sqrt u\), and the covolume of the coefficient integers
contributes the square root of the discriminant per rank.
Taking the two averages with the same place multiplicities
gives
\[
 \mu\le 2\lambda_s+2\log B+\log\operatorname{rd}(K_c)
              \le x^{2a+4}b .
\]
For the last estimate \(2q+1\le3r\), \(6f+48\le12f\) and
\(frL\le X\) give
\[
 \lambda_s\le760X^2 b.
\]

The other two contributions together are bounded by
\[
 1000X^2+217X^2=1217X^2.
\]
The factor \(x^4\) absorbs their sum
and the factor two on \(\lambda_s\) for \(x\ge10\).
\label{cor:parent-cusp-slope-4}

\proofstep{The actual ambient arithmetic lattice.}
On the actual row the ambient integral forms are the
\(\mathcal O_K\)-span of this lattice under the independent
semilinear realization. The norm of every geometric vector
at an extension \(\sigma\) is its geometric norm multiplied
by the positive scalar \(c_\sigma\) of the measure comparison
above. The direct and inverse comparisons amount to
\[
 e^{-122U}\le c_\sigma\le e^{122U}.
\]

Divide out this scalar. The normalized Gram on geometric
vectors extends to their complex span, by polarization and
density of the image of the totally complex field. Choosing
the conjugate embedding conjugates the Gram and preserves
its determinant. For the bound on arbitrary sublattices, let \(K'/K_c\) be a
normal extension containing \(K\), and write \(M_c\) for the
geometric lattice over \(K_c\). Give \(M_c\otimes\mathcal O_{K'}\)
the normalized metric. The Galois group permutes its places
by isometries, hence preserves the maximal destabilizing
saturated lattice. Its rational span descends to \(V_c\) over
\(K_c\).

To identify its integral lattice, tensor the exact
kernel sequence
\[
 0\longrightarrow M_c\cap V_c\longrightarrow M_c
       \longrightarrow (M_c\otimes K_c)/V_c
\]
with \(\mathcal O_{K'}\) over \(\mathcal O_{K_c}\).
Flatness identifies its kernel with the saturated
intersection over \(K'\). Normalized slopes of these extended
lattices are unchanged by the ramification and archimedean
multiplicity identities. Any sublattice over \(K\) also extends
at its original normalized slope. The maximal destabilizing
lattice over \(K'\) therefore bounds the slope over \(K\).
Restoring \(c_\sigma\) changes any squared slope by at most
\(244U\). Since \(U=fmbL\log\log x\le x^{a+1}b\) on our
family by \(m\le rn_V\) and \(\log\log x\le x\), we have
\(\mu_{\max}(\mathrm{lat}_{\rm all})\le x^A b\) for a fixed
\(A\), uniformly in the row field and the lift.
\label{cor:parent-cusp-slope-5}

\end{proof}

\paragraph{Rank on the two sides of the projection.}
The metric comparison leaves a rank factor on each side. To identify these ranks, use
the central projection to the retained tuples of
semisimple labels in the same finite group at the base level.
In the rational form frame the type space is the intersection
of the kernels of the character equations. Adding the kernel
of the retained projector gives the discarded space. Flatness
and the fact that the full integral forms span compute the two
lattice ranks from exactly these linear equations after any
complex embedding. Twisting by the inverse central form
character makes the first system the Whittaker character
equations. The group order clears the projection on the full
lattice, since its character coefficients times the order are
algebraic integers. Inversion conjugates these coefficients at
each embedding and the group acts unitarily, so the two parts
are orthogonal. \label{cor:parent-cusp-slope-6}

\begin{lemma}[Rank discarded in the primitive base frame]
\label{lem:parent-projector-rank-loss}
Use the retained central projector in
Proposition~\ref{prop:equivariant-selected-row}, transported to the
primitive base frame, with the finite label bounds of
Proposition~\ref{prop:paired-good-rank}. Put \(\eta=\exp(-x^{e_z/7})\).
The retained and discarded ranks \(d_g,d_b\) in the actual type lattice satisfy
\[
 d_b\le3\eta(d_b+d_g),\qquad d_g=D_0.
\]

\end{lemma}

\begin{proof}
\proofstep{Computing the finite bad-rank fraction.}
In the definite model the two groups \(H_k\) in
\eqref{eq:row-bad-share} act freely on the marked components.
On a nonterminal deep cover use the inverse images of \(P\)
(the character subgroup) and of \(H_k\) in its level
automorphism group. The pullback action on forms differs
from translation by a single character. Thus the product of
a translated separating section and a translated form still
has this character on the form factor. The invariant-sum map
in the proof of \eqref{eq:row-bad-share} embeds the restriction
to each inverse image \(H\) in \(a_H\) regular representations.

Let \(r_+,r_-\) be the ranks at weights \(\kappa+2c_0\) and \(\kappa-c_0\). The determinant injections and the rank growth at shift two on this cover give
\[
 a_H|H|\le r_+\le\tfrac32 r_-\le\tfrac32|\pi^{-1}P|\,d_{\rm type}.
\]
Here \(c_0\) is the weight
of the orbit-separating section on that cover, with the group
determinant nonzero on every component. The regular trace on
the inverse image reads the local semisimple-label share
downstairs, since the averaging factor for the kernel cancels
in the trace ratio. Using \eqref{eq:row-bad-types} for both
primes, each discarded share is at most
\(\eta=\exp(-x^{e_z/7})\) on the late family. Hence
\[
   d_b\le (3/2)(\eta+\eta)d_{\rm type}
       =3\eta(d_b+d_g).
\]
The definite regular comparison has factor one and satisfies
the same inequality. The retained image is the good space in
the primitive base frame, of rank \(d_g=D_0\).
\label{cor:parent-cusp-slope-7}
\end{proof}


\paragraph{The normalized covolume in the division bound.}
We can now express the covolume in the units used by the division estimate. At lift
\(j\) put
\[
 M_j=\sum_\rho(p_j-q_j)w_{j\rho}\log\Omega,
 \qquad U=B_0\log\log x ,
\]
where \((p_j,q_j)=(k+j-1,m-j)\) is the signature of the
base realization at weight \(\kappa_j\).

\begin{proposition}[Covolume on the selected base row]\label{prop:selected-base-covolume}
Fix the completed family of Proposition~\ref{prop:completed-division-pair}.
At nonterminal lifts take the model and level of
Lemma~\ref{hyp:covolume-deep-component} with the weight and
slope estimates of Proposition~\ref{prop:covolume-trace-descent};
at the terminal lift use Lemmas~\ref{lem:terminal-valuative-realization} and~\ref{lem:terminal-cm-type-slope}. Use the metric and
integral descent of Lemmas~\ref{lem:covolume-type-metric}--\ref{lem:covolume-integral-descent}.
There is a constant \(C_0\) such that the
normalized log covolume \(L_{3,j}\) of the good base lift of
rank \(D_0>0\) satisfies
\[
 L_{3,j}\le M_j+(C_0+1)U .
\]
The constant and scale cutoff can be taken before choosing the
translate, sequence, or lift.
\end{proposition}

\begin{proof}
\proofstep{The ambient type slope.}
Write \(C_0=244+\max(C_D,2(C_F+C_U)+9)\), where \(C_D\)
is the terminal constant. The common metric change costs at
most \(244U\) by Lemma~\ref{lem:covolume-metric-factor}. The two input
budgets for the coefficient of \(U\) are therefore
\[
 244+C_D\le C_0,\qquad 244+2(C_F+C_U)+9\le C_0 .
\]
They give \(\mu_{\rm type}\ge-M_i-C_0 U\) by the terminal type estimate or by
Proposition~\ref{prop:covolume-trace-descent} with the metric
comparison of Lemmas~\ref{lem:covolume-type-metric} and~\ref{lem:covolume-metric-factor}. In the nonterminal case the heights and positivity
used in the wall calculation are those of Sections~\ref{subsec:flag-point-heights}
and~\ref{subsec:flag-wall-positivity} on this component. Their constants
enter the vanishing wall error at fixed model. The smaller-block tail
estimate controls the change of field and weight.

Let \(e_g\) be the self-adjoint good idempotent and
\(e_b=1-e_g\). The order \(s_0\) of the projection group clears
their denominators and satisfies \(2\log s_0\le b\). For the
integral type lattice \(\Lambda\) put
\(\Lambda_a=\Lambda\cap e_a\Lambda_{\mathbb Q}\),
\(d_a=\operatorname{rk}\Lambda_a\) for \(a=g,b\).
Projection gives an injection
\[
 \Lambda/(\Lambda_g\oplus\Lambda_b)
 \lhook\joinrel\longrightarrow s_0^{-1}\Lambda_b/\Lambda_b,
 \qquad [v]\longmapsto[e_bv].
\]
Indeed \(s_0e_bv\in\Lambda_b\); if \(e_bv\in\Lambda_b\), then
\(v-e_bv\in\Lambda_g\). The normalized squared-degree index
is at most \(2d_b\log s_0\). Orthogonality at all infinite
places makes the determinant metric multiplicative, hence
\[
 (d_g+d_b)\mu_{\rm type}\le
 d_g\mu_g+d_b(\mu_{\max}+2\log s_0),
\]
with \(\mu_{\max}\) the maximal slope of the ambient
lattice.

The preceding ambient slope and regular trace estimates give
\(d_b\le3\eta(d_g+d_b)\) with
\(\eta=e^{-x^{e_z/7}}\) and
\(\mu_{\max}\le x^A b\) for a fixed exponent \(A\).
Take \(B'\) such that \(|M_i|\le x^{B'}b\).
For \(\eta\le1/12\), \(d_b\le4\eta d_g\). Put
\(Q=(x^A+x^{B'}+C_0+1)U\), an upper bound for each of
\(\mu_{\max},2\log s_0,-\mu_{\rm type}\). Rearranging the degree
inequality gives
\[
 \mu_g\ge\mu_{\rm type}
 -\frac{d_b}{d_g}(\mu_{\max}+2\log s_0-\mu_{\rm type})
 \ge\mu_{\rm type}-12\eta Q .
\]
Since \(e_z>0\) is fixed, \(12\eta Q\le U\) eventually uniformly.
\label{cor:covolume-component-transfer-22}

Transport the good lattice into the independent base frame,
whose rank is \(D_0=d_g>0\). The normalized norm at each
embedding equals the geometric norm times \(s=1/\sqrt{n_P}\le1\),
with \(n_P\) the number of indices in the level normalization.
Log covolume per rank in the squared metric therefore satisfies
\[
 L_{3,i}=-\mu_g+\operatorname{mean}_{v\mid\infty}\log s^2
              \le M_i+(C_0+1)U .
\]
This is the base-level estimate used before the sharp-level
comparison at the same row. The arithmetic application in
Section~\ref{sec:word-division} requires both this estimate and the direct-row functional bound.
If their thresholds on the fixed family are \(x_{\rm cov}\) and
\(x_{\rm fun}\), use
\[
 x\ge\max\{x_{\rm cov},x_{\rm fun}\}.
\]
Both estimates then hold for every translate, sequence and lift
allowed by Proposition~\ref{prop:selected-base-covolume}.
\label{cor:covolume-component-transfer-23}

The two blocks used in the Hodge calculation are precisely
the \(p_j\) positive and \(q_j\) negative entries of the
base realization at weight \(\kappa_j\). The flag Chern calculation
uses these two block ranks, and the determinant shift and theta-root
twist use the weights of this base realization.
In particular the negative block enters the covolume through
the signed signature. The geometric model and the form
frame keep the cut, weight and complex Gram when identifying
the base row, so the signed term in the fixed-scale determinant
is precisely \(M_j\). At each orientation the auxiliary
evaluation lattice used for the four-term index has its own
finite chart conditions. These are checked on the common cusp
space before passage to the base and sharp lattices.
\label{cor:covolume-component-transfer-24}
\end{proof}
